\paragraph{角色投影、平滑塌缩与外置算术核的函子分裂}
在模剩余 Fourier 反演、window-\(6\) 群商障碍、末位充分化、输出位势的解析腔室刚性、二步原子的 primitive/finite-part 位移、零温四态地基层分裂以及局部化数系的 Pontryagin 对偶链已经建立之后，还可把离散波纹与连续平滑之间的关系进一步压缩为下列七条彼此闭合的结论。它们分别刻画非平凡角色扇区、非自由群商障碍、末位二点终端塌缩、单解析相内的波滑分裂、两步原子的 Abel 有限部位移、零温平滑包络的有限状态化，以及连续相位层之后仍不可消去的 profinite 素数核。

\begin{theorem}[角色投影的扭曲扇区分裂]\label{thm:xi-time-part67-character-projection-twisted-sectors}
\leanverified{paper\_xi\_time\_part67\_character\_projection\_twisted\_sectors}
固定整数 \(q\ge 2\)。设 \(B_{q,a}(n)\) 为定义在 \(a\in \ZZ/(q\ZZ)\) 上的计数函数，并令
\[
\widehat B_{q,j}(n):=\sum_{a\in \ZZ/(q\ZZ)}B_{q,a}(n)\,\omega_q^{ja},
\qquad
\omega_q:=e^{2\pi i/q}.
\]
则有 Fourier 反演
\[
B_{q,a}(n)=\frac1q\sum_{j=0}^{q-1}\widehat B_{q,j}(n)\,\omega_q^{-ja},
\qquad
\frac1q\sum_{b\in \ZZ/(q\ZZ)}B_{q,b}(n)=\frac1q\,\widehat B_{q,0}(n),
\]
因而中心化波纹恰为
\[
B_{q,a}(n)-\frac1q\,\widehat B_{q,0}(n)
=
\frac1q\sum_{j=1}^{q-1}\widehat B_{q,j}(n)\,\omega_q^{-ja}.
\]
对真实输入 \(40\) 状态核取 \(B_{q,a}(n)=A_{q,a}(n)\) 或 \(B_{q,a}(n)=P_{q,a}(n)\)，则
\[
\widehat A_{q,j}(n)=a_n(\omega_q^j),
\qquad
\widehat P_{q,j}(n)=p_n(\omega_q^j),
\]
从而平滑均值恰是平凡角色扇区，而全部离散波纹严格由非平凡单位根扭曲 \(M(\omega_q^j)\) 承载。进一步，若记
\[
\rho_q:=\max_{1\le j\le q-1}\rho\!\bigl(M(\omega_q^j)\bigr),
\qquad
\theta_q:=\frac{\log\rho_q}{\log\lambda},
\qquad
\lambda=\varphi^2,
\]
则这些非平凡扇区在 periodic 与 primitive 两层均以 \(O(\rho_q^n)\) 衰减；并且存在明确模数 \(q\) 使 \(\theta_q>\frac12\)，且该性质沿倍数塔向上遗传。
\end{theorem}
\begin{proof}
有限循环群 \(\ZZ/(q\ZZ)\) 上的离散 Fourier 反演直接给出前两式；中心化公式只需把 \(j=0\) 项单独提出。

对真实输入 \(40\) 状态核，定理 \ref{thm:xi-dirichlet-residue-fourier-parseval-exact} 已在 periodic 与 primitive 两层分别给出
\[
a_n(\omega_q^j)=\sum_{a\in \ZZ/(q\ZZ)}A_{q,a}(n)\,\omega_q^{ja},
\qquad
p_n(\omega_q^j)=\sum_{a\in \ZZ/(q\ZZ)}P_{q,a}(n)\,\omega_q^{ja},
\]
故扭曲 trace 正是 residue counts 的非平凡角色分量。命题 \ref{prop:real-input-40-output-dirichlet-nonrh} 进一步给出：这些扭曲分量在 periodic/primitive 两层都以 \(O(\rho_q^n)\) 衰减；并且存在明确模数使 \(\theta_q>\frac12\)，其中 \(q=2\) 已足够。最后，由结论 \ref{con:xi-time-part9l-dirichlet-twist-divisibility-monotonicity}，一旦某个模数满足 \(\theta_q>\frac12\)，则其全部倍数也都满足同一严格不等式。证毕。
\end{proof}

\begin{theorem}[window-\texorpdfstring{$6$}{6} 平滑包络的非自由群商障碍]\label{thm:xi-time-part67-window6-nonfree-group-quotient-obstruction}
设
\[
\Omega_6:=\{0,1\}^6,
\qquad
F_6:=\Fold_6^{\mathrm{bin}}:\Omega_6\to X_6.
\]
若试图把 \(F_6\) 的纤维压缩解释为某个固定有限群的轨道商，或解释为某个阿贝尔线性陪集划分，则该方案必失败。更具体地，\(F_6\) 的纤维直方图
\[
2:8,\qquad 3:4,\qquad 4:9
\]
迫使任何与之同轨道的群作用都必须带三层稳定子分层；而在保持 fold-gauge 纤维独立置换语义时，完备异常签名仍需
\[
\Aut(\Fold_6^{\mathrm{bin}})^{\mathrm{ab}}\cong (\ZZ/2\ZZ)^{21}.
\]
因此，审计窗中的平滑包络所消去的并非单一统一对称冗余，而只能是随纤维变化的稳定子残差；其自然承载对象是带变稳定子的群胚型商，而非单一群商。
\end{theorem}
\begin{proof}
定理 \ref{thm:xi-window6-group-quotient-isotropy-strata} 已证明：任何把 \(F_6\) 纤维分解实现为群轨道分解的方案，都不可能是自由作用，也不可能退化为等势轨道；并且轨道大小只能取 \(2,3,4\)，从而稳定子恰分成三层。于是，固定自由群商与阿贝尔线性陪集解释同时被排除。

另一方面，定理 \ref{thm:xi-window6-oracle-anomaly-coordinate-bifurcation} 证明：若继续保持 fold-gauge 的纤维独立置换语义，则完备异常信息的最小承载对象仍为
\(
(\ZZ/2\ZZ)^{21}
\)；
任何把异常坐标压到 \(21\) 以下的方案都必须引入跨纤维关系，从而破坏单一纤维直积语义。故这里被“平滑”掉的并不是某个固定群的统一冗余，而是随纤维变化的稳定子残差；用群胚而非单一群来描述，正是这一变稳定子机制的自然形式化。证毕。
\end{proof}

\begin{theorem}[末位投影的渐近充分性与终端二点塌缩]\label{thm:xi-time-part67-lastbit-terminal-csiszar-collapse}
设
\[
\ell_m(w):=w_m,
\qquad
\mu_m(w):=\frac{d_m^{\mathrm{bin}}(w)}{2^m},
\qquad
u_m(w):=\frac{1}{|X_m|}
\qquad (w\in X_m).
\]
则 \(\ell_m\) 对分布对 \((\mu_m,u_m)\) 是渐近充分统计量：对任意凸函数 \(f:(0,\infty)\to\RR\) 满足 \(f(1)=0\)，有
\[
D_f(\mu_m\Vert u_m)
=
D_f\!\bigl((\ell_m)_\#\mu_m\Vert (\ell_m)_\#u_m\bigr)
+O_f\!\Bigl(\Bigl(\frac{\varphi}{2}\Bigr)^m\Bigr).
\]
并且若记
\[
p_*:=\frac{1}{\varphi\sqrt5},
\qquad
q_*:=\varphi^{-2},
\]
则
\[
\lim_{m\to\infty}D_f(\mu_m\Vert u_m)
=
D_f\!\bigl(\mathrm{Bernoulli}(p_*)\Vert \mathrm{Bernoulli}(q_*)\bigr)
=
\frac1\varphi\,f\!\left(\frac{\varphi^2}{\sqrt5}\right)
+
\frac1{\varphi^2}\,f\!\left(\frac{\varphi}{\sqrt5}\right).
\]
换言之，相对于稳定层均匀基线的全部渐近 Csisz\'ar 几何，都在末位二点实验上终端塌缩。
\end{theorem}
\begin{proof}
定理 \ref{thm:xi-fold-lastbit-asymptotic-sufficiency} 已在全变差与 Le Cam 意义下给出末位投影 \(\ell_m\) 的渐近充分性；下面只需把这一充分化提升到 \(f\)-散度层面。

由定理 \ref{thm:xi-time-part64b-fold-information-geometry-two-point-collapse}，
\[
\frac{\mu_m(w)}{u_m(w)}
=
\frac{\varphi^{2-\ell_m(w)}}{\sqrt5}
+O\!\Bigl(\Bigl(\frac{\varphi}{2}\Bigr)^m\Bigr)
\qquad (w\in X_m)
\]
在 \(w\) 上一致成立。于是存在紧区间 \(K\subset(0,\infty)\)，使得对充分大的 \(m\)，上式右端与左端都落在 \(K\) 内。由于凸函数在紧子区间上局部 Lipschitz，故对每个固定 \(f\) 有
\[
D_f(\mu_m\Vert u_m)
=
\sum_{b\in\{0,1\}}u_m(\ell_m=b)\,
f\!\left(\frac{\varphi^{2-b}}{\sqrt5}\right)
+O_f\!\Bigl(\Bigl(\frac{\varphi}{2}\Bigr)^m\Bigr).
\]

另一方面，设
\[
p_m:=((\ell_m)_\#\mu_m)(\{1\}),
\qquad
q_m:=((\ell_m)_\#u_m)(\{1\})=\frac{F_m}{F_{m+2}}.
\]
对上式在每个末位块上取 \(u_m\)-平均，得到
\[
\frac{p_m}{q_m}
=
\frac{\varphi}{\sqrt5}
+O\!\Bigl(\Bigl(\frac{\varphi}{2}\Bigr)^m\Bigr),
\qquad
\frac{1-p_m}{1-q_m}
=
\frac{\varphi^2}{\sqrt5}
+O\!\Bigl(\Bigl(\frac{\varphi}{2}\Bigr)^m\Bigr),
\]
从而
\[
D_f\!\bigl((\ell_m)_\#\mu_m\Vert (\ell_m)_\#u_m\bigr)
=
q_m\,f\!\left(\frac{p_m}{q_m}\right)
+(1-q_m)\,f\!\left(\frac{1-p_m}{1-q_m}\right)
\]
与前一式只差 \(O_f((\varphi/2)^m)\)。这便证明了渐近充分性。

最后，由
\[
q_m=\frac{F_m}{F_{m+2}}\longrightarrow \frac1{\varphi^2},
\qquad
1-q_m=\frac{F_{m+1}}{F_{m+2}}\longrightarrow \frac1{\varphi},
\]
以及上面的两个比值极限，立即得到
\[
\lim_{m\to\infty}D_f\!\bigl((\ell_m)_\#\mu_m\Vert (\ell_m)_\#u_m\bigr)
=
\frac1\varphi\,f\!\left(\frac{\varphi^2}{\sqrt5}\right)
+\frac1{\varphi^2}\,f\!\left(\frac{\varphi}{\sqrt5}\right).
\]
再与刚证得的误差估计合并，即得所示极限公式。证毕。
\end{proof}

\begin{theorem}[单一解析腔室内的波滑分裂]\label{thm:xi-time-part67-single-analytic-chamber-wave-smooth-splitting}
在真实输入 \(40\) 状态核的输出位势参数化
\[
u=e^\theta
\]
下，审计窗口
\[
|\theta|<\theta_0,
\qquad
\theta_0=0.0533477827\ldots,
\]
满足
\[
u\in (e^{-\theta_0},e^{\theta_0})
\Subset
(0,u_R),
\qquad
u_R\approx 3.59262181344.
\]
因而该窗口既不跨越正实退化点，也完全停留在
\[
\Lambda(u)<\lambda_+(u)^{1/2}
\]
的统一亚临界区间内。对于任何可归约为解析 \(\arg\max\) 读出的折叠谱观测量，若其在此窗口内出现离散波纹与连续平滑的分裂，则该分裂必发生于同一解析相内部，而不能归因于谱分支碰撞或正实相变。
\end{theorem}
\begin{proof}
推论 \ref{cor:xi-collision-u1-audit-window-uniform-subcriticality} 已证明：
\[
|\theta|<\theta_0
\Longrightarrow
u=e^\theta\in(0.932837\ldots,1.054796\ldots)\Subset(0,u_R),
\]
从而整个 \(u=1\) 邻域审计窗都处在统一亚临界腔室内，不会接触平方根门槛 \(\Lambda(u)=\lambda_+(u)^{1/2}\) 的临界端点。

另一方面，命题 \ref{prop:xi-time-part9m3-analytic-chamber-rigidity-sampling-readout} 说明：在任何不穿越退化点的实解析区间上，凡由解析 \(\arg\max\) 读出定义的纤维计数、容量与矩谱，都只能在离散交叉集之外分段常值。于是，在本窗口内若仍观察到“波纹--平滑”二分，则其来源只能是同一解析腔室内部不同读出函子的分裂，而不能被解释为谱支交换或正实相变。证毕。
\end{proof}

\begin{theorem}[两步原子剥离与 Abel 有限部位移]\label{thm:xi-time-part67-two-step-atomic-peeling-finitepart-shift}
对真实输入 \(40\) 状态核，有精确分解
\[
\Delta(z,u)=(1-uz^2)\,F(z,u),
\qquad
P(z,u)=uz^2+P_{\mathrm{core}}(z,u).
\]
在未加权情形 \(u=1\) 下，若 \(z_\star=\varphi^{-2}\)，并记 \(\mathfrak M,\mathfrak M_{\mathrm{core}}\) 分别为 full/core 的 Abel 有限部常数，则
\[
\log \mathfrak M
=
\log \mathfrak M_{\mathrm{core}}+\varphi^{-4},
\qquad
\varphi^{-4}=\frac{7-3\sqrt5}{2}.
\]
因此，由离散波纹统计转入平滑统计并非中性平均，而是一次精确的两步 primitive Euler 原子剥离。
\end{theorem}
\begin{proof}
定理 \ref{thm:xi-time-part65-atomic-primitive-abel-shift} 已直接给出
\[
P(z,u)=uz^2+P_{\mathrm{core}}(z,u),
\qquad
\log\mathfrak M=\log\mathfrak M_{\mathrm{core}}+\varphi^{-4}.
\]
其中 \(\varphi^{-4}=(7-3\sqrt5)/2\) 是同一定理中的闭式常数。于是 full/core 之差在 primitive 层只作用于长度 \(2\) 的单一原子，而在 Abel 有限部层只留下一个完全闭式的代数位移。证毕。
\end{proof}

\begin{theorem}[零温平滑包络的有限状态化]\label{thm:xi-time-part67-zero-temp-envelope-finite-stateization}
对真实输入 \(40\) 状态核的输出位势 \(P(\theta)=\log\lambda(e^\theta)\)，有
\[
P'(\theta)\longrightarrow \frac12
\qquad (\theta\to+\infty),
\]
并且满足平方根端点律
\[
\frac12-P'(\theta)=O(e^{-\theta/2}).
\]
另一方面，在零温缩放
\[
z=\sqrt{w/u}
\]
下，
\[
\lim_{u\to+\infty}\Delta\!\left(\sqrt{w/u},u\right)
=
(1-w)\det(I-wB_\infty),
\]
其中
\[
B_\infty=
\begin{pmatrix}
0&1&1&0\\
0&0&1&0\\
0&0&3&1\\
1&0&0&3
\end{pmatrix}.
\]
因此，零温极限中被选出的平滑包络并非整个 \(40\) 状态核的粗粒化投影，而是由显式四状态地基 SFT 所控制的有限状态对象；原子层仅留下因子 \((1-w)\)，而主导极点由地基层而非原子层决定。
\end{theorem}
\begin{proof}
定理 \ref{thm:xi-output-potential-zero-temp-square-root-critical-law} 给出
\[
\frac12-P'(\theta)=\frac{d}{2c}e^{-\theta/2}+O(e^{-\theta}),
\]
故零温端点的饱和斜率确为 \(1/2\)。

再由定理 \ref{thm:xi-zero-temp-atomic-ground-exact-splitting}，
\[
\lim_{u\to+\infty}\Delta\!\left(\sqrt{w/u},u\right)
=(1-w)\det(I-wB_\infty),
\]
这说明零温解析对象严格分裂为一个 primitive 原子因子与一个四状态地基 SFT 的动力学行列式。最后，推论 \ref{cor:xi-zero-temp-leading-pole-ground-dominance} 进一步表明正实最内层零点来自 \(\det(I-wB_\infty)\) 而非 \((1-w)\)，故主导平滑尺度由四状态地基层决定。证毕。
\end{proof}

\begin{theorem}[连续相位层与外置算术核]\label{thm:xi-time-part67-continuous-phase-external-arithmetic-kernel}
设 \(\varnothing\neq S\subset\mathbb P\) 为有限素数集，并记
\[
G_S:=\ZZ[S^{-1}].
\]
则其 Pontryagin 对偶满足短正合列
\[
0\longrightarrow \prod_{p\in S}\ZZ_p
\longrightarrow \widehat{G_S}
\longrightarrow \TT
\longrightarrow 0,
\]
并且该 profinite 核 \(\prod_{p\in S}\ZZ_p\) 可由相位采样函数 \(\Sigma^{\mathrm{ph}}_{G_S}\) 唯一恢复。因而圆相位因子 \(\TT\) 只是连通可见层，真正的算术内容则以外置 prime-register 核的形式保留在 profinite 方向上。

更进一步，不存在任何有限秩加法账本能够把正整数乘法幺半群忠实严格同态化；因此，上述 profinite 核不能被有限多个加法寄存器吸收。
\end{theorem}
\begin{proof}
定理 \ref{thm:xi-localized-phase-sampling-closed-form} 已证明：\(\Sigma^{\mathrm{ph}}_{G_S}\) 唯一恢复素数集 \(S\)，并因而唯一恢复对偶短正合列中的 profinite 核
\[
\prod_{p\in S}\ZZ_p.
\]
定理 \ref{thm:xi-localized-basic-extension-strictly-nonsplit} 进一步说明该短正合列在 \(S\neq\varnothing\) 时严格不分裂，因此 \(\TT\) 不能替代整个对偶对象；它只给出连通商，而非全部算术信息。

最后，定理 \ref{thm:xi-time-part64b-local-additivity-nonlocal-multiplicativity-polarization} 已排除一切把
\(
(\NN_{\times},\times)
\)
严格忠实同态化到有限秩加法账本的方案。故这里的 prime-register profinite 核确实是不可由有限加法寄存器消化的外置算术部分。证毕。
\end{proof}

\noindent
上述七条结果表明，离散间隔统计中的波纹并不应被理解为连续对象自身的粗糙化。在本文框架内，它们恰是非平凡角色扇区、window-\(6\) 的变稳定子残差、两步 primitive 原子与 prime-register profinite 核在离散读出层的共同显影；相应地，平滑极限则由平凡角色投影、末位二点塌缩、Abel 有限部反项剥离以及零温地基层选择共同构成。因而，“波”与“滑”并非同一对象的两种视觉描画，而是同一算术动力学对象在不同函子层上的两种范畴像。

\paragraph{谱临界、输出慢模、端点有限部与熵缺损常数的统一闭合}
在碰撞位势的解析半径、输出位势的四阶累积结构、地基层端点能垒、两步 primitive 原子、极端偏置有限部以及 bin-fold 的信息位率阈值都已建立之后，还可把这些分散结论进一步压缩为下列七条彼此闭合的终端后果。它们分别刻画碰撞 RH 临界与解析分歧的严格分离、输出模扭曲慢模的四阶刚性、端点代价与地基熵的精确配平、输出大偏差相图的双坐标代数化、七条分支二回路在 primitive 层的单原子塌缩、极端偏置有限部对单一四次代数参数的依赖，以及熵缺损常数在三种信息尺度上的同一化。

\begin{theorem}[碰撞 RH 临界先于解析分歧出现]\label{thm:xi-time-part68-collision-rh-critical-before-analytic-branch}
设
\[
P_2(\theta):=\log \rho(e^\theta),
\]
并记 \(R_\theta\) 为其在 \(\theta=0\) 处 Taylor 级数的收敛半径。又设 \(u_R>1\) 为核版 RH/Ramanujan 条件
\[
\Lambda(u)\le \lambda_1(u)^{1/2}
\]
失效的唯一正临界点，\(\theta_R:=\log u_R\)。则
\[
0<\theta_R<R_\theta.
\]
更具体地，
\[
u_R\approx 3.59262181344,\qquad
\theta_R\approx 1.27888,\qquad
R_\theta\approx 2.76230289346087.
\]
因此，碰撞核的 RH/Ramanujan 失效并非由解析延拓的首个分歧点触发，而是在解析域内部由实谱不等式
\[
\Lambda(u)\le \lambda_1(u)^{1/2}
\]
的翻转所触发；该临界属于纯谱型临界，而非分支型临界。
\end{theorem}
\begin{proof}
定理 \ref{thm:xi-rh-critical-before-nearest-complex-branch} 已证明：\(P_2\) 的全部代数分歧点恰由判别式方程
\[
4u^3+25u^2+88u+8=0
\]
的三根经
\[
\theta=\log u_j+2\pi i k
\qquad (k\in\ZZ)
\]
产生，且最近分歧点的模长正是
\[
R_\theta\approx 2.76230289346087.
\]
同一定理又给出：\(u_R\) 是五次方程
\[
u^5+4u^4+3u^3-96u^2-42u-14=0
\]
的唯一正根，\(\theta_R=\log u_R\approx 1.27888\)，并且
\[
\textsf{RH}(u)\Longleftrightarrow 0<u\le u_R.
\]
比较数值便得
\[
0<\theta_R<R_\theta.
\]
于是对任意 \(\theta\in(\theta_R,R_\theta)\)，一方面 \(|\theta|<R_\theta\)，故 \(P_2(\theta)\) 仍位于 \(\theta=0\) 的解析收敛盘内；另一方面 \(u=e^\theta>u_R\)，故 RH/Ramanujan 条件已经失效。再由同一定理中的
\[
\Lambda(u_R)=\lambda_1(u_R)^{1/2},
\]
可知临界边界确由谱不等式的等号翻转刻画，而非由解析分支奇性触发。证毕。
\end{proof}

\begin{theorem}[输出模扭曲慢模的四阶累积刚性]\label{thm:xi-time-part68-output-twist-slowmode-fourth-cumulant-rigidity}
记
\[
\lambda(u):=\rho(M(u)),
\qquad
\lambda:=\lambda(1)=\varphi^2,
\qquad
P(\theta):=\log \lambda(e^\theta),
\qquad
\kappa_2:=P''(0),
\qquad
\kappa_4:=P^{(4)}(0).
\]
则
\[
\kappa_2,\kappa_4\in\QQ(\sqrt5),
\qquad
\kappa_2>0.
\]
对主单位根扭曲通道设
\[
t_m:=\frac{2\pi}{m},
\qquad
\rho_{m,1}:=\rho\!\bigl(M(e^{it_m})\bigr).
\]
则有严格渐近展开
\[
\log \frac{\rho_{m,1}}{\lambda}
=
-\frac{\kappa_2}{2}\Bigl(\frac{2\pi}{m}\Bigr)^2
+\frac{\kappa_4}{24}\Bigl(\frac{2\pi}{m}\Bigr)^4
+O(m^{-6}).
\]
并且在已审计模数 \(m=5,10,20\) 上，四阶预测的绝对误差分别约为
\[
1.08\times 10^{-6},\qquad
2.78\times 10^{-8},\qquad
6.21\times 10^{-10},
\]
而仅保留二阶项时的相应误差分别约为
\[
5.67\times 10^{-4},\qquad
3.66\times 10^{-5},\qquad
2.30\times 10^{-6}.
\]
因此，输出模扭曲的近平凡慢模在四阶以内已被前四阶累积量刚性锁定；四阶项并非可忽略修饰，而是恢复慢模几何所必需的主导校正。
\end{theorem}
\begin{proof}
命题 \ref{prop:real-input-40-output-cumulants-closed} 已给出
\[
P^{(k)}(0)\in\QQ(\sqrt5),
\]
特别是 \(\kappa_2,\kappa_4\in\QQ(\sqrt5)\) 且 \(\kappa_2>0\)。再由定理 \ref{thm:xi-output-potential-dirichlet-fixed-j-fourth-order} 在 \(j=1\) 时的结论，
\[
\log \frac{\rho_{m,1}}{\lambda}
=
-\frac{\kappa_2}{2}\Bigl(\frac{2\pi}{m}\Bigr)^2
+\frac{\kappa_4}{24}\Bigl(\frac{2\pi}{m}\Bigr)^4
+O(m^{-6}),
\]
这就给出所示四阶展开。

另一方面，附录注 \ref{rem:real-input-40-dirichlet-cumulant-error-order} 已逐项记录 \(m=5,10,20\) 上的四阶预测误差与纯二阶高斯近似误差，恰为题述数值。再由定理 \ref{thm:xi-output-potential-dirichlet-fourth-order-gain}，
\[
\frac{\widetilde\varepsilon_m^{(2)}}{\widetilde\varepsilon_m^{(4)}}=\Omega(m^2),
\]
故四阶修正不仅改善常数，而且在阶数量级上严格压缩误差。于是四阶项确为慢模重建所必需的主导校正。证毕。
\end{proof}

\begin{proposition}[端点代价缺口恰等于地基熵]\label{prop:xi-time-part68-endpoint-gap-ground-entropy-identity}
\leanverified{paper\_xi\_time\_part68\_endpoint\_gap\_ground\_entropy\_identity}
设 \(I(\alpha)\) 为输出位势
\[
P(\theta)=\log\lambda(e^\theta)
\]
的归一化大偏差速率函数，并记
\[
\alpha_{\max}:=\lim_{\theta\to+\infty}P'(\theta)=\frac12,
\qquad
h_{\mathrm{ground}}:=\log c.
\]
则有精确恒等式
\[
I(0)-I(\alpha_{\max})=h_{\mathrm{ground}}.
\]
更具体地，
\[
I(0)=\log(\varphi^2),
\qquad
I\!\left(\frac12\right)=\log\!\left(\frac{\varphi^2}{c}\right),
\qquad
h_{\mathrm{ground}}=\log c.
\]
因此，右端点的大偏差代价相对于左端点的全部减免，恰由地基层的组合熵指数精确补偿。
\end{proposition}
\begin{proof}
定理 \ref{thm:xi-output-potential-endpoint-gap-ground-state} 已直接给出
\[
I(0)-I\!\left(\frac12\right)=h_{\mathrm{ground}}.
\]
而推论 \ref{cor:real-input-40-alpha-max} 说明
\[
\alpha_{\max}=\frac12,
\]
推论 \ref{cor:real-input-40-ground-entropy} 则给出
\[
h_{\mathrm{ground}}=\log c.
\]
再结合推论 \ref{cor:real-input-40-rate-endpoints} 中的端点闭式
\[
I(0)=\log(\varphi^2),
\qquad
I(1/2)=\log(\varphi^2/c),
\]
即可得到题述恒等式与具体数值表达。证毕。
\end{proof}

\begin{theorem}[输出大偏差相图的代数一致化]\label{thm:xi-time-part68-output-ldp-phase-diagram-algebraization}
设 \(u=e^\theta\)，\(\lambda(u)\) 为方程
\[
G(\lambda,u)=0
\]
的 Perron 正根，并定义平衡态下输出 \(1\) 的典型密度
\[
\alpha(u):=\frac{\mathrm d}{\mathrm d\theta}\log\lambda(e^\theta).
\]
则存在非零多项式
\[
F(\alpha,u)\in\QQ[\alpha,u],
\qquad
H(\alpha,\lambda)\in\QQ[\alpha,\lambda]
\]
使得沿 Perron 分支恒有
\[
F(\alpha(u),u)=0,
\qquad
H(\alpha(u),\lambda(u))=0.
\]
并且归一化大偏差速率函数满足参数公式
\[
I(\alpha(u))
=
\alpha(u)\log u-\log\frac{\lambda(u)}{\lambda(1)}.
\]
因此，输出大偏差相图在 \((\alpha,u)\) 与 \((\alpha,\lambda)\) 两组坐标中都可代数化；真正的超越性只保留在 Legendre 对偶中的单一对数项。
\end{theorem}
\begin{proof}
命题 \ref{prop:real-input-40-pressure-freezing} 已给出输出位势的特征多项式
\[
G(\lambda,u)\in\QQ[\lambda,u],
\]
并说明 \(\lambda(u)\) 是其 Perron 正根。另一方面，附录 \ref{app:real-input-40-zeta-u} 中的输出位势大偏差代数参数化已把 \(\alpha(u)\) 写成
\[
\alpha(u)=\left.\frac{N(\lambda,u)}{D(\lambda,u)}\right|_{\lambda=\lambda(u)},
\qquad
N,D\in\QQ[\lambda,u].
\]
因此沿 Perron 分支有
\[
\alpha\,D(\lambda,u)-N(\lambda,u)=0,
\qquad
G(\lambda,u)=0.
\]
对 \(\lambda\) 取结果式并约去共同因子，便得到某个非零多项式
\[
F(\alpha,u)\in\QQ[\alpha,u]
\]
满足 \(F(\alpha(u),u)=0\)。同理，对 \(u\) 取结果式并约去共同因子，可得某个非零多项式
\[
H(\alpha,\lambda)\in\QQ[\alpha,\lambda]
\]
满足 \(H(\alpha(u),\lambda(u))=0\)。

最后，由同一附录中的参数公式，或者直接由严格凸压力的驻点 Legendre 公式，
\[
I(\alpha(u))
=
\alpha(u)\log u-\log\frac{\lambda(u)}{\lambda(1)}.
\]
于是 \((\alpha,u)\) 与 \((\alpha,\lambda)\) 两种坐标下的相图都由代数关系控制，而超越部分只出现在该单一对数项中。证毕。
\end{proof}

\begin{theorem}[七条分支二回路在 primitive 层塌缩为唯一原子]\label{thm:xi-time-part68-seven-branch-two-cycles-collapse-single-atom}
设输出位势的 primitive 生成函数为
\[
P(z,u)=\sum_{n\ge 1}p_n(u)z^n.
\]
则在真实输入 \(40\) 状态核上有精确分解
\[
P(z,u)=uz^2+P_{\mathrm{core}}(z,u).
\]
另一方面，branch 层审计表明，满足
\[
(|\tau|,E)=(2,1)
\]
的两步回路恰有 \(7\) 条。因而，这 \(7\) 条几何上可区分的 branch 二回路在 primitive--Witt 化之后并不留下 \(7\) 个独立原子，而是完全塌缩为唯一长度 \(2\) 原子
\[
uz^2.
\]
\end{theorem}
\begin{proof}
定理 \ref{thm:xi-time-part62d-atomic-two-step-witt-rankone-collapse} 已直接给出
\[
P(z,u)=uz^2+P_{\mathrm{core}}(z,u),
\]
并同时指出 branch 层中满足
\[
(|\tau|,E)=(2,1)
\]
的两步回路总数恰为 \(7\)。逐项比较系数可知，full/core 的 primitive 差分只支撑在长度 \(2\) 的单一坐标上；等价地，
\[
\dim_{\QQ(u)}
\operatorname{span}\{\,p_n(u)-p_n^{\mathrm{core}}(u):n\ge 1\,\}=1.
\]
因此，这七条 branch 两步回路在 primitive--Witt 层只能留下唯一原子 \(uz^2\)。证毕。
\end{proof}

\begin{theorem}[极端偏置有限部由单一四次代数参数完全控制]\label{thm:xi-time-part68-extreme-bias-finitepart-single-quartic-control}
记
\[
b:=b_\infty:=\rho(B_\infty)^{-1}\in(0,1),
\]
等价地，\(b\) 是四次方程
\[
1-6b+9b^2-b^3-b^4=0
\]
在区间 \((0,1)\) 内的最小正根。则极端偏置 \(u\to+\infty\) 下，主极点常数 \(C_\infty\) 与 Abel 有限部常数 \(\mathfrak M_\infty\) 满足
\[
C_\infty=
\frac{1}{2b(1-b)(6-18b+3b^2+4b^3)},
\]
\[
\log\mathfrak M_\infty
=
\log C_\infty
-\sum_{m\ge 2}\frac{\mu(m)}{m}
\log\Bigl[(1-b^m)\bigl(1-6b^m+9b^{2m}-b^{3m}-b^{4m}\bigr)\Bigr].
\]
数值上
\[
b\approx 0.27807480214113,\qquad
C_\infty\approx 1.89745149363,\qquad
\log\mathfrak M_\infty\approx 0.282290255262.
\]
因此，极端偏置下的全部有限部数据都由单一四次代数参数 \(b\) 完整统辖。
\end{theorem}
\begin{proof}
命题 \ref{prop:real-input-40-logM-infty} 已直接证明：令
\[
b=\lim_{u\to\infty}u\,z_\star(u)^2=\frac{1}{c^2},
\]
则 \(b\) 满足题述四次方程，并且 \(C_\infty\) 与 \(\log\mathfrak M_\infty\) 的闭式恰为上式。另一方面，推论 \ref{cor:xi-zero-temp-leading-pole-ground-dominance} 说明
\[
b=b_\infty=\rho(B_\infty)^{-1}
\]
正是零温缩放极限中控制主导正实零点的位置参数。故无论从主极点常数还是从 Abel 有限部常数看，极端偏置下的主导有限部数据都确实由同一个四次代数参数 \(b\) 完整决定。证毕。
\end{proof}

\begin{theorem}[熵缺损常数的三重同一化]\label{thm:xi-time-part68-entropy-deficit-triple-identification}
记
\[
\mu_m(w):=\frac{d_m^{\mathrm{bin}}(w)}{2^m},
\qquad
W_m\sim\mu_m,
\qquad
\gamma_m:=\frac{1}{m}\log d_m^{\mathrm{bin}}(W_m).
\]
则对任意 \(q>0\)（其中 \(q=1\) 按 Shannon 熵解释）都有
\[
\lim_{m\to\infty}\frac{1}{m}\bigl(m\log 2-H_q(\mu_m)\bigr)
=
\log\frac{2}{\varphi},
\qquad
\gamma_m\longrightarrow \log\frac{2}{\varphi}.
\]
若一列二进制预算 \(B_m\in\NN\) 满足
\[
\mathrm{Succ}_m^{\mathrm{bin}}(B_m)\ge 1-\varepsilon
\qquad
(0<\varepsilon<1\ \text{固定}),
\]
则必有
\[
\liminf_{m\to\infty}\frac{B_m}{m}
\ge
\log_2\frac{2}{\varphi}.
\]
因此，同一个常数 \(\log(2/\varphi)\) 同时控制全部正阶 Rényi 熵的广延缺损、典型微态的纤维多重度指数以及保持固定成功率所需的最小外置比特率。
\end{theorem}
\begin{proof}
由定理 \ref{thm:fold-bin-renyi-rate-collapse}，
\[
\frac{1}{m}H_q(\mu_m)\longrightarrow \log\varphi
\qquad (q>0),
\]
并且
\[
\gamma_m\longrightarrow \log\frac{2}{\varphi}.
\]
于是
\[
\frac{1}{m}\bigl(m\log 2-H_q(\mu_m)\bigr)
=
\log 2-\frac{1}{m}H_q(\mu_m)
\longrightarrow
\log 2-\log\varphi
=
\log\frac{2}{\varphi},
\]
这就得到前两条极限。

另一方面，若用 \(B_m\) 个二进制外置信息位实现恢复，则其状态数为 \(M_m=2^{B_m}\)。推论 \ref{cor:xi-reconstruction-cardinality-strong-converse} 对 \(\delta=1-\varepsilon\) 给出
\[
B_m
\ge
m-\log_2F_{m+2}+\log_2(1-\varepsilon).
\]
两边同除以 \(m\)，再用 Binet 渐近
\[
\log_2F_{m+2}=m\log_2\varphi+O(1),
\]
便得
\[
\liminf_{m\to\infty}\frac{B_m}{m}\ge \log_2\frac{2}{\varphi}.
\]
故题述三种量确由同一个常数统一控制。证毕。
\end{proof}

\noindent
上述七条结果表明，碰撞核的 RH 破裂、输出单位根慢模、地基层端点能垒、两步 primitive 原子、零温有限部常数与 bin-fold 的熵缺损阈值并不是彼此分离的局部现象。它们共同揭示：一方面，碰撞侧的临界首先发生在解析盘内部，因而属于纯谱型翻转；另一方面，输出侧的慢模、端点代价与大偏差相图都受同一条代数压力曲线及其四阶累积结构支配；最终，这种谱--代数刚性又在 primitive 原子剥离、零温有限部与信息预算下界中外显为同一组可审计的不变量。因而，从碰撞 RH 窗口到输出大偏差，再到部分可逆预算与极端偏置有限部，全文中的主导常数与临界门槛实际上共享同一条谱--组合--信息几何闭合链。

\paragraph{低预算线性相、二点实验塌缩、局部化周期谱与形式线性化的闭合后果}
在审计偶窗的最小退化 Fibonacci 律、末位两层塌缩、局部化 solenoid 的自覆盖分类、素数寄存器的幂等估值算术化以及 fold-gauge 缺陷的纯奇次重整化已经建立之后，还可把这些链条继续压缩为下列七条结果。它们分别刻画审计偶窗中的低预算严格线性段及其首个偏离阈值、bin-fold 输出实验的 Le Cam 二点塌缩、一般局部化有理乘法自同态的固定点与 Artin--Mazur \(\zeta\) 函数、相应周期点几何的指数率、素数寄存器幂等元的精确闭式计数，以及 fold-gauge 缺陷的形式 Koenigs 线性化。

\begin{theorem}[bin-fold 审计偶数窗上的低预算严格线性律]\label{thm:xi-time-part69-audited-even-pure-pigeonhole-linear-phase}
对审计偶数窗
\[
m\in\{6,8,10,12\},
\]
定义
\[
d_{\min}(m):=\min_{w\in X_m}d_m^{\mathrm{bin}}(w),
\qquad
C_m^{\mathrm{bin}}(B):=\sum_{w\in X_m}\min\bigl(d_m^{\mathrm{bin}}(w),2^B\bigr),
\]
\[
\mathrm{Succ}_m^{\mathrm{bin}}(B):=\frac{C_m^{\mathrm{bin}}(B)}{2^m}.
\]
若整数预算 \(B\ge 0\) 满足
\[
2^B\le d_{\min}(m)=F_{m/2},
\]
则有精确公式
\[
C_m^{\mathrm{bin}}(B)=2^B\,F_{m+2},
\qquad
\mathrm{Succ}_m^{\mathrm{bin}}(B)=2^{B-m}F_{m+2}.
\]
因此，在最小退化阈值之前，bin-fold 的每增加一比特辅助信息都会使可逆微态总数严格翻倍，且不出现任何次级修正。
\end{theorem}
\begin{proof}
定理 \ref{thm:fold-bin-min-degeneracy-fib-audited} 已证明：当
\[
m\in\{6,8,10,12\}
\]
时，
\[
d_{\min}(m)=F_{m/2}.
\]
若 \(2^B\le F_{m/2}\)，则对所有 \(w\in X_m\) 都有
\[
2^B\le d_{\min}(m)\le d_m^{\mathrm{bin}}(w),
\]
从而
\[
\min\bigl(d_m^{\mathrm{bin}}(w),2^B\bigr)=2^B.
\]
因此
\[
C_m^{\mathrm{bin}}(B)
=
\sum_{w\in X_m}2^B
=
2^B|X_m|.
\]
另一方面，由稳定词计数律
\[
|X_m|=F_{m+2},
\]
即得
\[
C_m^{\mathrm{bin}}(B)=2^BF_{m+2}.
\]
再除以 \(2^m\)，得到
\[
\mathrm{Succ}_m^{\mathrm{bin}}(B)=2^{B-m}F_{m+2}.
\]
证毕。
\end{proof}

\begin{corollary}[bin-fold 容量曲线的首个折点恰位于 \texorpdfstring{$F_{m/2}$}{F(m/2)}]\label{cor:xi-time-part69-binfold-first-kink}
在定理 \ref{thm:xi-time-part69-audited-even-pure-pigeonhole-linear-phase} 的设定下，容量曲线
\[
B\longmapsto C_m^{\mathrm{bin}}(B)
\]
的首个偏离线性区间发生在最小整数
\[
B_{\mathrm{kink}}(m):=\min\{B\in\NN:\ 2^B>F_{m/2}\}
=
\bigl\lfloor \log_2 F_{m/2}\bigr\rfloor+1.
\]
更准确地，
\[
C_m^{\mathrm{bin}}(B)=2^BF_{m+2}
\quad\Longleftrightarrow\quad
2^B\le F_{m/2},
\]
而一旦 \(2^B>F_{m/2}\)，便有严格不等式
\[
C_m^{\mathrm{bin}}(B)<2^BF_{m+2}.
\]
\end{corollary}
\begin{proof}
“若”方向即定理 \ref{thm:xi-time-part69-audited-even-pure-pigeonhole-linear-phase}。反之，由 \(X_m\) 的有限性与
\[
d_{\min}(m)=F_{m/2}
\]
可取 \(w_0\in X_m\) 使
\[
d_m^{\mathrm{bin}}(w_0)=F_{m/2}.
\]
若 \(2^B>F_{m/2}\)，则
\[
\min\bigl(d_m^{\mathrm{bin}}(w_0),2^B\bigr)=F_{m/2}<2^B,
\]
而其余各项至多为 \(2^B\)。于是
\[
C_m^{\mathrm{bin}}(B)<2^B|X_m|=2^BF_{m+2}.
\]
故首个折点恰由阈值 \(F_{m/2}\) 决定。证毕。
\end{proof}

\begin{theorem}[bin-fold 统计实验的 Le Cam 二点塌缩]\label{thm:xi-time-part69-binfold-lecam-two-point-collapse}
\leanverified{paper\_xi\_time\_part69\_binfold\_lecam\_two\_point\_collapse}
设
\[
\mu_m(w):=\frac{d_m^{\mathrm{bin}}(w)}{2^m},
\qquad
u_m(w):=\frac{1}{|X_m|},
\qquad
\ell_m(w):=w_m\in\{0,1\},
\]
并定义两块
\[
A_0:=\{w\in X_m:\ w_m=0\},
\qquad
A_1:=\{w\in X_m:\ w_m=1\}.
\]
令
\[
T_m:=\ell_m:X_m\to\{0,1\},
\qquad
K_m(i,\cdot):=\mathrm{Unif}(A_i)\qquad(i=0,1).
\]
再定义极限二点实验
\[
Q_\infty:=\left(\frac1\varphi,\frac1{\varphi^2}\right),
\qquad
P_\infty:=\left(\frac{\varphi}{\sqrt5},\frac{1}{\varphi\sqrt5}\right),
\]
其中坐标顺序对应 \((0,1)\)。则原始两点统计实验
\[
\mathcal E_m:=\{u_m,\mu_m\}\quad\text{定义在 }X_m
\]
与二元实验
\[
\mathcal E_\infty:=\{Q_\infty,P_\infty\}\quad\text{定义在 }\{0,1\}
\]
在 Le Cam 意义下渐近等价。更具体地，
\[
D_{\mathrm{TV}}\!\bigl(T_{m\#}u_m,Q_\infty\bigr)=O(\varphi^{-2m}),
\qquad
D_{\mathrm{TV}}\!\bigl(T_{m\#}\mu_m,P_\infty\bigr)=O\!\Bigl(\Bigl(\frac{\varphi}{2}\Bigr)^m\Bigr),
\]
并且
\[
D_{\mathrm{TV}}\!\bigl(u_m,K_{m\#}Q_\infty\bigr)=O(\varphi^{-2m}),
\qquad
D_{\mathrm{TV}}\!\bigl(\mu_m,K_{m\#}P_\infty\bigr)=O\!\Bigl(\Bigl(\frac{\varphi}{2}\Bigr)^m\Bigr).
\]
因此
\[
\Delta(\mathcal E_m,\mathcal E_\infty)\xrightarrow[m\to\infty]{}0.
\]
\end{theorem}
\begin{proof}
先看均匀基线 \(u_m\)。由
\[
|A_0|=F_{m+1},
\qquad
|A_1|=F_m,
\qquad
|X_m|=F_{m+2},
\]
可得
\[
T_{m\#}u_m=
\left(\frac{F_{m+1}}{F_{m+2}},\frac{F_m}{F_{m+2}}\right)
\]
且由 Binet 公式比较相邻 Fibonacci 比值可知，其两坐标分别与 \(Q_\infty\) 的对应坐标相差 \(O(\varphi^{-2m})\)。特别地，
\[
D_{\mathrm{TV}}\!\bigl(T_{m\#}u_m,Q_\infty\bigr)=O(\varphi^{-2m}).
\]
另一方面，\(u_m\) 在每个块 \(A_i\) 上本来就是严格均匀分布，因此
\[
u_m=K_{m\#}(T_{m\#}u_m).
\]
由 Markov 核的收缩性，
\[
D_{\mathrm{TV}}\!\bigl(u_m,K_{m\#}Q_\infty\bigr)
\le
D_{\mathrm{TV}}\!\bigl(T_{m\#}u_m,Q_\infty\bigr)
=
O(\varphi^{-2m}).
\]

再看 \(\mu_m\)。定理 \ref{thm:xi-fold-lastbit-asymptotic-sufficiency} 已给出
\[
D_{\mathrm{TV}}\!\bigl(\mu_m,K_{m\#}(T_{m\#}\mu_m)\bigr)
=
O\!\Bigl(\Bigl(\frac{\varphi}{2}\Bigr)^m\Bigr).
\]
另一方面，定理 \ref{thm:xi-fold-lastbit-two-layer-collapse} 证明 \(\mu_m\) 与末位两层模型 \(\nu_m(w)=\varphi^{-m-w_m}\) 在全变差意义下相差
\[
O\!\Bigl(\Bigl(\frac{\varphi}{2}\Bigr)^m\Bigr).
\]
而由
\[
\nu_m(A_0)=\varphi^{-m}F_{m+1}=\frac{\varphi}{\sqrt5}+O(\varphi^{-2m}),
\qquad
\nu_m(A_1)=\varphi^{-m-1}F_m=\frac{1}{\varphi\sqrt5}+O(\varphi^{-2m}),
\]
再与
\[
D_{\mathrm{TV}}(\mu_m,\nu_m)=O\!\Bigl(\Bigl(\frac{\varphi}{2}\Bigr)^m\Bigr)
\]
合并，可知
\[
D_{\mathrm{TV}}\!\bigl(T_{m\#}\mu_m,P_\infty\bigr)
=
O\!\Bigl(\Bigl(\frac{\varphi}{2}\Bigr)^m\Bigr).
\]
于是再用 Markov 核收缩性，
\[
\begin{aligned}
D_{\mathrm{TV}}\!\bigl(\mu_m,K_{m\#}P_\infty\bigr)
&\le
D_{\mathrm{TV}}\!\bigl(\mu_m,K_{m\#}(T_{m\#}\mu_m)\bigr)
+
D_{\mathrm{TV}}\!\bigl(K_{m\#}(T_{m\#}\mu_m),K_{m\#}P_\infty\bigr)\\
&\le
O\!\Bigl(\Bigl(\frac{\varphi}{2}\Bigr)^m\Bigr)
+
D_{\mathrm{TV}}\!\bigl(T_{m\#}\mu_m,P_\infty\bigr)\\
&=
O\!\Bigl(\Bigl(\frac{\varphi}{2}\Bigr)^m\Bigr).
\end{aligned}
\]

最后，对参数集 \(\{0,1\}\) 上的两点实验，前向缺陷由核 \(T_m\) 控制，后向缺陷由核 \(K_m\) 控制，因此 Le Cam 距离满足
\[
\Delta(\mathcal E_m,\mathcal E_\infty)
\le
\max\Bigl\{
D_{\mathrm{TV}}(T_{m\#}u_m,Q_\infty),
D_{\mathrm{TV}}(T_{m\#}\mu_m,P_\infty),
D_{\mathrm{TV}}(u_m,K_{m\#}Q_\infty),
D_{\mathrm{TV}}(\mu_m,K_{m\#}P_\infty)
\Bigr\},
\]
右端趋于 \(0\)，故结论成立。证毕。
\end{proof}

\begin{theorem}[局部化 solenoid 上有理乘法自同态的固定点与 Artin--Mazur \texorpdfstring{$\zeta$}{zeta} 函数]\label{thm:xi-time-part69-localized-rational-endomorphism-fixedpoints-zeta}
\leanverified{paper\_xi\_time\_part69\_localized\_rational\_endomorphism\_fixedpoints\_zeta}
设 \(S\subset\mathbb P\) 为有限素数集，并记
\[
G_S:=\ZZ[S^{-1}],
\qquad
\Sigma_S:=\widehat{G_S}.
\]
取有理数
\[
r=\frac{a}{b}\in G_S,
\qquad
\gcd(a,b)=1,
\]
其中 \(b\) 的全部素因子都属于 \(S\)，并设 \(f_r:\Sigma_S\to\Sigma_S\) 为乘法映射
\[
x\longmapsto rx\qquad (x\in G_S)
\]
的 Pontryagin 对偶连续自同态。若 \(|a|\neq |b|\)，则对每个 \(n\ge 1\)，都有
\[
\#\mathrm{Fix}(f_r^n)=|a^n-b^n|_{S^\perp},
\]
其中对任意非零整数 \(N\)，记
\[
|N|_{S^\perp}:=(|N|)_{S^\perp}
\]
为其 \(S\)-自由部分。并且 Artin--Mazur 动力学 \(\zeta\) 函数满足
\[
\zeta_{f_r}(z)
:=
\exp\!\left(
\sum_{n\ge 1}\frac{\#\mathrm{Fix}(f_r^n)}{n}z^n
\right)
=
\exp\!\left(
\sum_{n\ge 1}\frac{|a^n-b^n|_{S^\perp}}{n}z^n
\right).
\]
换言之，局部化 solenoid 上的周期点几何只感知 \(a^n-b^n\) 的 \(S\)-外部。
\end{theorem}
\begin{proof}
由定理 \ref{thm:xi-localized-solenoid-self-cover-degree-classification}，任意连续自同态都唯一对应于某个 \(r\in G_S\)，其对偶离散同态为
\[
x\longmapsto rx
\qquad (x\in G_S).
\]
因此对每个 \(n\ge 1\)，
\[
f_r^n
\]
在对偶离散侧对应于乘法
\[
x\longmapsto r^n x.
\]
按定理 \ref{thm:xi-localized-solenoid-fixedpoints-artin-mazur-series} 的同样论证，
\[
\mathrm{Fix}(f_r^n)\cong \widehat{G_S/(r^n-1)G_S}.
\]

现把
\[
r^n-1=\frac{a^n-b^n}{b^n}
\]
代入。由于 \(b\) 的全部素因子都属于 \(S\)，故 \(b^n\in G_S^\times\) 是单位，从而
\[
(r^n-1)G_S=(a^n-b^n)G_S.
\]
又因 \(|a|\neq |b|\)，故 \(a^n-b^n\neq 0\)。于是由定理 \ref{thm:xi-localized-quotient-exact-sfree-externalization}，
\[
G_S/(a^n-b^n)G_S
\cong
\ZZ/|a^n-b^n|_{S^\perp}\ZZ.
\]
取 Pontryagin 对偶便得到
\[
\#\mathrm{Fix}(f_r^n)=|a^n-b^n|_{S^\perp}.
\]
再将该固定点计数代入 Artin--Mazur \(\zeta\) 函数定义，即得所示级数公式。证毕。
\end{proof}

\begin{theorem}[局部化有理自覆盖的周期点指数率只由分子分母的最大模控制]\label{thm:xi-time-part69-localized-rational-endomorphism-growth-rate}
沿用定理 \ref{thm:xi-time-part69-localized-rational-endomorphism-fixedpoints-zeta} 的记号，并进一步假设
\[
|a|\neq |b|.
\]
则有极限
\[
\lim_{n\to\infty}\frac{1}{n}\log \#\mathrm{Fix}(f_r^n)
=
\log\max\{|a|,|b|\}.
\]
换言之，局部化 solenoid 上有理自覆盖的周期点指数率不受 \(S\)-内素因子的影响，而仅由分子与分母模长中的较大者决定。
\end{theorem}
\begin{proof}
由定理 \ref{thm:xi-time-part69-localized-rational-endomorphism-fixedpoints-zeta} 证明中的同一对偶化论证，对每个 \(n\ge 1\) 都有
\[
\mathrm{Fix}(f_r^n)\cong \widehat{G_S/(r^n-1)G_S}.
\]
又因 \(|a|\neq |b|\)，故 \(a^n-b^n\neq 0\)。并且
\[
r^n-1=\frac{a^n-b^n}{b^n},
\]
其中 \(b^n\in G_S^\times\) 为单位，于是
\[
(r^n-1)G_S=(a^n-b^n)G_S.
\]
再由定理 \ref{thm:xi-localized-quotient-exact-sfree-externalization}，
\[
\#\mathrm{Fix}(f_r^n)=|a^n-b^n|_{S^\perp}.
\]
因此只需比较 \(a^n-b^n\) 的 \(S\)-内与 \(S\)-外部分的增长。

先固定 \(p\in S\)。若 \(p\mid ab\)，则由 \(\gcd(a,b)=1\) 可知 \(p\) 不能同时整除 \(a\) 与 \(b\)，从而对任意 \(n\ge 1\) 都有
\[
p\nmid (a^n-b^n).
\]
若 \(p\nmid ab\)，则 \(a/b\) 在 \((\ZZ/p\ZZ)^\times\) 中有定义。记其模 \(p\) 的乘法阶为 \(t_p\)。只有在 \(t_p\mid n\) 时，\(p\) 才可能整除 \(a^n-b^n\)；并且由标准 lifting-the-exponent 估计，存在仅依赖于 \(p,a,b\) 的常数 \(C_p\)，使得
\[
v_p(a^n-b^n)\le C_p+v_p(n)
\qquad(n\ge 1).
\]
由于 \(S\) 有限，遂得
\[
\sum_{p\in S}v_p(a^n-b^n)\log p=O(\log n),
\]
因而
\[
\frac{1}{n}\sum_{p\in S}v_p(a^n-b^n)\log p\longrightarrow 0.
\]

另一方面，若 \(|a|>|b|\)，则
\[
|a^n-b^n|
=
|a|^n\Bigl|1-(b/a)^n\Bigr|
=
|a|^n(1+o(1)),
\]
故
\[
\frac{1}{n}\log |a^n-b^n|\longrightarrow \log|a|.
\]
若 \(|b|>|a|\)，同理得到极限为 \(\log|b|\)。

最后，由
\[
\log |a^n-b^n|_{S^\perp}
=
\log |a^n-b^n|
-\sum_{p\in S}v_p(a^n-b^n)\log p
\]
即知
\[
\frac{1}{n}\log |a^n-b^n|_{S^\perp}
\longrightarrow
\log\max\{|a|,|b|\}.
\]
再代回
\(\#\mathrm{Fix}(f_r^n)=|a^n-b^n|_{S^\perp}\)，即得所示极限。证毕。
\end{proof}

\begin{theorem}[素数寄存器半群幂等元的精确计数律]\label{thm:xi-time-part69-prime-register-idempotent-count-law}
沿用定理 \ref{thm:killo-finite-semigroup-prime-valuation-arithmetization} 的记号，设
\[
S_n:=\left\{\prod_{i=0}^{n-1}q_i^{a_i}: a_i\in\{0,\dots,n-1\}\right\},
\]
并记
\[
E_n:=\{N\in S_n:\ N\star N=N\},
\qquad
E_{n,k}:=\{N\in E_n:\ |\operatorname{Im}(f_N)|=k\}
\qquad (1\le k\le n).
\]
则有精确计数公式
\[
|E_{n,k}|=\binom{n}{k}k^{n-k}
\qquad (1\le k\le n),
\]
从而
\[
|E_n|=\sum_{k=1}^{n}\binom{n}{k}k^{n-k}.
\]
换言之，素数寄存器中的幂等投影按像集大小完全分层，而其总数与 \([n]\) 上幂等自映射的经典计数严格一致。
\end{theorem}
\begin{proof}
由定理 \ref{thm:killo-finite-semigroup-prime-valuation-arithmetization} 与推论 \ref{cor:killo-projection-idempotent-prime-register}，\(N\in S_n\) 为幂等元当且仅当对应映射
\[
f_N:[n]\to[n]
\]
为幂等变换。

固定 \(k\in\{1,\dots,n\}\)，并设 \(|\operatorname{Im}(f_N)|=k\)。记
\[
A:=\operatorname{Im}(f_N)\subseteq [n],
\qquad |A|=k.
\]
幂等条件
\[
f_N\circ f_N=f_N
\]
等价于：每个 \(a\in A\) 都必须是定点，即
\[
f_N(a)=a
\qquad (a\in A),
\]
而对每个 \(x\in[n]\setminus A\)，其像只需任意落在 \(A\) 中。因而，一旦像集 \(A\) 固定，幂等变换数恰为
\[
k^{n-k}.
\]
再对全部 \(\binom{n}{k}\) 个 \(k\)-元像集求和，即得
\[
|E_{n,k}|=\binom{n}{k}k^{n-k}.
\]
最后对 \(k=1,\dots,n\) 求和，得到
\[
|E_n|=\sum_{k=1}^{n}\binom{n}{k}k^{n-k}.
\]
证毕。
\end{proof}

\begin{theorem}[fold-gauge 缺陷的形式 Koenigs 线性化]\label{thm:xi-time-part69-fold-gauge-formal-koenigs-linearization}
\leanverified{paper\_xi\_time\_part69\_fold\_gauge\_formal\_koenigs\_linearization}
记
\[
E_m:=\log C_m-\log(2\pi),
\qquad
x:=\frac{\varphi}{2}\in(0,1).
\]
由定理 \ref{thm:xi-time-part60ab-gauge-constant-recursive-bernoulli-stripping}，存在奇形式级数
\[
F(y):=\sum_{r\ge 1}a_r y^{2r-1}\in y\RR[[y^2]],
\qquad
a_1=\frac{1}{3\varphi}\neq 0,
\]
使
\[
E_m=F(x^m).
\]
则存在唯一奇形式幂级数
\[
\mathcal K(z)=3\varphi\,z+\kappa_3 z^3+\kappa_5 z^5+\cdots \in z\RR[[z^2]]
\]
使得对任意 \(N\ge 1\)，若记 \(\mathcal K_N\) 为其截断到 \(z^{2N+1}\) 的多项式，则
\[
\mathcal K_N(E_m)=x^m+O\!\bigl(x^{(2N+1)m}\bigr)
\qquad (m\to\infty).
\]
等价地，若 \(\mathcal R(z)\in z\RR[[z^2]]\) 为定理 \ref{thm:xi-time-part64a-fold-gauge-purely-odd-renormalization} 中的唯一形式重整化映射，则
\[
\mathcal K(\mathcal R(z))=x\,\mathcal K(z)
\]
在形式幂级数意义下成立。其首个非平凡系数为
\[
\kappa_3=\frac{135+63\sqrt5}{40}.
\]
\end{theorem}
\begin{proof}
由于 \(a_1\neq 0\)，级数 \(F\) 在形式幂级数环中存在唯一复合逆
\[
F^{-1}(z)\in z\RR[[z]].
\]
又因为 \(F\) 为奇级数，其逆也必为奇级数，于是
\[
\mathcal K(z):=F^{-1}(z)\in z\RR[[z^2]].
\]
由构造立得
\[
\mathcal K(E_m)=\mathcal K(F(x^m))=x^m
\]
作为形式恒等式成立。若 \(\mathcal K_N\) 为截断到 \(z^{2N+1}\) 的多项式，则
\[
\mathcal K(z)-\mathcal K_N(z)=O(z^{2N+3}),
\]
而 \(E_m\sim a_1x^m\)，故
\[
\mathcal K_N(E_m)=x^m+O(x^{(2N+3)m}),
\]
特别地也有较弱的
\[
\mathcal K_N(E_m)=x^m+O(x^{(2N+1)m}).
\]

另一方面，定理 \ref{thm:xi-time-part64a-fold-gauge-purely-odd-renormalization} 已给出唯一形式重整化映射
\[
\mathcal R(z)=F\!\bigl(xF^{-1}(z)\bigr).
\]
于是
\[
\mathcal K(\mathcal R(z))
=
F^{-1}\!\bigl(F(xF^{-1}(z))\bigr)
=
xF^{-1}(z)
=
x\,\mathcal K(z),
\]
这正是 Koenigs 线性化方程。

最后计算首项系数。由 \(r=1,2\) 的系数公式，
\[
a_1=\frac{1}{3\varphi},
\qquad
a_2=\frac{B_4}{2\cdot 3}\,(\varphi^{-1}+\varphi)
=
-\frac{\sqrt5}{180},
\]
其中用到 \(B_4=-1/30\) 与 \(\varphi+\varphi^{-1}=\sqrt5\)。设
\[
\mathcal K(z)=k_1 z+k_3 z^3+O(z^5).
\]
把它代入
\[
\mathcal K(F(y))=y
\]
并比较 \(y\) 与 \(y^3\) 的系数，得到
\[
k_1a_1=1,
\qquad
k_1a_2+k_3a_1^3=0.
\]
第一式给出
\[
k_1=\frac{1}{a_1}=3\varphi.
\]
第二式给出
\[
k_3=-\frac{a_2}{a_1^4}
=
\frac{135+63\sqrt5}{40}.
\]
故 \(\kappa_3=k_3\)。证毕。
\end{proof}

\noindent
上述七条结果表明，审计偶窗的低预算容量、其首个非线性折点、bin-fold 的完整统计实验、局部化 solenoid 的周期点几何、素数寄存器的有限幂等结构以及 fold-gauge 缺陷的形式重整化之间存在同一类低维闭合机制：在有限分辨率层面，它表现为严格线性段、精确折点与精确的幂等计数；在统计层面，它表现为末位二点实验对全模型的 Le Cam 塌缩；在局部化动力学层面，它既表现为 \(S\)-外部因子对周期点计数的完全控制，也表现为指数率仅由 \(\max\{|a|,|b|\}\) 锁定；而在形式动力学层面，它最终收束为 Koenigs 坐标下的严格乘法线性化。于是，有限预算盲区、容量折点、二点统计极限、局部化周期谱与重整化正规形不再是彼此分离的技术现象，而被统一压缩为同一条可审计的降维主线。

\paragraph{单位自同构、周期指数与 prime-register 盲性}
在一般局部化有理自同态的固定点与指数率已经获得闭式之后，单位参数区间还可进一步收束为一条更强的动力学链：连续自同构群完全由单位格给出，周期点公式退化为 \(S\)-自由差分的精确读数，而遍历、mixing 与同构分类之间则出现了严格分离。

\begin{theorem}[局部化 solenoid 的连续自同构群完全分类]\label{thm:xi-time-part69-localized-unit-automorphism-group-classification}
\leanverified{paper\_xi\_time\_part69\_localized\_unit\_automorphism\_group\_classification}
设 \(S\subset\mathbb P\) 为有限素数集，并记
\[
G_S:=\ZZ[S^{-1}],
\qquad
\Sigma_S:=\widehat{G_S}.
\]
则连续群自同构群满足显式同构
\[
\boxed{
\operatorname{Aut}_{\mathrm{cts}}(\Sigma_S)
\cong
G_S^\times
=
\left\{\pm\prod_{p\in S}p^{n_p}:n_p\in\ZZ\right\}
\cong
C_2\times \ZZ^{|S|}.
}
\]
\end{theorem}
\begin{proof}
定理 \ref{thm:xi-localized-solenoid-self-cover-degree-classification} 已给出
\[
\operatorname{Aut}_{\mathrm{cts}}(\Sigma_S)\cong G_S^\times
=
\left\{\pm\prod_{p\in S}p^{n_p}:n_p\in\ZZ\right\}.
\]
由于 \(S\) 有限，映射
\[
\varepsilon\prod_{p\in S}p^{n_p}
\longmapsto
\bigl(\varepsilon,(n_p)_{p\in S}\bigr),
\qquad
\varepsilon\in\{\pm1\},
\]
给出该单位群与 \(C_2\times \ZZ^{|S|}\) 的群同构，故结论成立。证毕。
\end{proof}

\begin{theorem}[单位比值自同构的周期点公式]\label{thm:xi-time-part69-unit-automorphism-periodic-point-formula}
\leanverified{paper\_xi\_time\_part69\_unit\_automorphism\_periodic\_point\_formula}
沿用定理 \ref{thm:xi-time-part69-localized-unit-automorphism-group-classification} 的记号。取
\[
q\in G_S^\times\setminus\{\pm1\},
\]
并将其写成既约形式
\[
q=\frac{a}{b},
\qquad
a\in\ZZ\setminus\{0\},
\qquad
b\in\NN,
\qquad
\gcd(a,b)=1,
\qquad
\operatorname{Supp}(|ab|)\subseteq S.
\]
令
\[
\alpha_q:\Sigma_S\to\Sigma_S
\]
为与离散端乘法 \(x\mapsto qx\) 对偶的连续群自同构。则对任意 \(m\ge 1\)，有
\[
\boxed{
\operatorname{Fix}(\alpha_q^m)
\cong
\ZZ/|a^m-b^m|_{S^\perp}\ZZ,
\qquad
\#\operatorname{Fix}(\alpha_q^m)=|a^m-b^m|_{S^\perp}.
}
\]
\end{theorem}
\begin{proof}
由定理 \ref{thm:xi-time-part69-localized-unit-automorphism-group-classification}，\(\alpha_q\) 确为 \(\Sigma_S\) 的连续群自同构。再把定理 \ref{thm:xi-time-part69-localized-rational-endomorphism-fixedpoints-zeta} 应用于 \(r=q=a/b\)，便得到
\[
\#\operatorname{Fix}(\alpha_q^m)=|a^m-b^m|_{S^\perp},
\]
以及对应的有限循环群结构。证毕。
\end{proof}

\begin{theorem}[单位自同构的周期指数与 prime-register 的多项式级盲性]\label{thm:xi-time-part69-unit-automorphism-polynomial-primeblindness}
\leanpartial{paper\_xi\_time\_part69\_unit\_automorphism\_polynomial\_primeblindness}{当前 Lean 目标按给定签名为假：取 $a=b=1$ 时假设成立而结论要求正的下界，但左侧恒为 $0$。}
沿用定理 \ref{thm:xi-time-part69-unit-automorphism-periodic-point-formula} 的记号，并设
\[
q\neq \pm1.
\]
则有极限
\[
\boxed{
\lim_{m\to\infty}\frac1m\log \#\operatorname{Fix}(\alpha_q^m)
=
\log\max\{|a|,b\}.
}
\]
更精确地，存在常数 \(C=C(q,S)>0\)，使得对充分大的 \(m\) 都有
\[
\boxed{
C^{-1}\,\frac{\max\{|a|,b\}^{\,m}}{m^{|S|}}
\le
\#\operatorname{Fix}(\alpha_q^m)
\le
C\,\max\{|a|,b\}^{\,m}.
}
\]
因此，prime-register \(S\) 只能改变周期点计数的多项式级修正，而不能改变其指数主率。
\end{theorem}
\begin{proof}
由定理 \ref{thm:xi-time-part69-unit-automorphism-periodic-point-formula}，
\[
\#\operatorname{Fix}(\alpha_q^m)
=
|a^m-b^m|_{S^\perp}
=
\frac{|a^m-b^m|}{\prod_{p\in S}p^{v_p(a^m-b^m)}}.
\]
记
\[
M:=\max\{|a|,b\},
\qquad
\lambda:=\frac{\min\{|a|,b\}}{M}\in(0,1).
\]
由于 \(q\neq \pm1\)，故 \(|a|\neq b\)。若 \(M=|a|\)，则
\[
|a^m-b^m|
=
M^m\left|\operatorname{sgn}(a)^m-\lambda^m\right|;
\]
若 \(M=b\)，则
\[
|a^m-b^m|
=
M^m\left|(\operatorname{sgn}(a)\lambda)^m-1\right|.
\]
两种情形都表明：存在常数 \(c_0,C_0>0\)，使得对充分大的 \(m\)，有
\[
c_0 M^m\le |a^m-b^m|\le C_0 M^m.
\]

现估计 \(S\)-内素因子。对固定 \(p\in S\)：

若 \(p\mid ab\)，则由 \(\gcd(a,b)=1\) 可知 \(p\) 只能整除 \(a,b\) 之一，从而对任意 \(m\ge 1\) 都有
\[
v_p(a^m-b^m)=0.
\]
在这种情形下取 \(C_p:=0\)。
若 \(p\nmid ab\)，则 \(a/b\) 在 \((\ZZ/p\ZZ)^\times\) 中有定义。记其模 \(p\) 的乘法阶为 \(d_p\)。当 \(d_p\nmid m\) 时，显然
\[
v_p(a^m-b^m)=0.
\]
当 \(d_p\mid m\) 时，由标准的 lifting-the-exponent 估计（在 \(p=2\) 时使用其 dyadic 变体）可知，存在仅依赖于 \(p,a,b\) 的常数 \(C_p\)，使得
\[
v_p(a^m-b^m)\le C_p+v_p(m).
\]
因此对所有 \(p\in S\) 与所有 \(m\ge1\)，统一有
\[
v_p(a^m-b^m)\le C_p+v_p(m).
\]
于是
\[
\prod_{p\in S}p^{v_p(a^m-b^m)}
\le
\left(\prod_{p\in S}p^{C_p}\right)\prod_{p\in S}p^{v_p(m)}
\le
C_1\,m^{|S|},
\]
其中 \(C_1:=\prod_{p\in S}p^{C_p}\)。

把上述估计与 \(c_0 M^m\le |a^m-b^m|\le C_0 M^m\) 合并，即得
\[
\#\operatorname{Fix}(\alpha_q^m)
=
\frac{|a^m-b^m|}{\prod_{p\in S}p^{v_p(a^m-b^m)}}
\ge
\frac{c_0}{C_1}\,\frac{M^m}{m^{|S|}}
\]
以及
\[
\#\operatorname{Fix}(\alpha_q^m)\le C_0 M^m.
\]
重新吸收常数便得到所示双边界。两边取对数、除以 \(m\) 并令 \(m\to\infty\)，立得指数极限。证毕。
\end{proof}

\begin{theorem}[rank-\texorpdfstring{$1$}{1} 局部化世界中遍历与 mixing 完全塌缩为 \texorpdfstring{$q\neq\pm1$}{q!=±1}]\label{thm:xi-time-part69-unit-automorphism-ergodic-mixing-collapse}
\leanverified{paper\_xi\_time\_part69\_unit\_automorphism\_ergodic\_mixing\_collapse}
沿用定理 \ref{thm:xi-time-part69-localized-unit-automorphism-group-classification} 的记号。对任意 \(q\in G_S^\times\)，记
\[
\alpha_q:\Sigma_S\to\Sigma_S
\]
为对偶于乘法 \(x\mapsto qx\) 的连续群自同构。则在 Haar 概率测度意义下，
\[
\boxed{
\alpha_q\ \text{遍历}
\iff
\alpha_q\ \text{mixing}
\iff
q\neq \pm1.
}
\]
\end{theorem}
\begin{proof}
记 \(\mu_S\) 为 \(\Sigma_S\) 上的 Haar 概率测度。对每个 \(x\in G_S\)，记对应字符为
\[
\chi_x(\omega):=\omega(x)
\qquad(\omega\in \Sigma_S).
\]
则 Pontryagin 对偶关系给出
\[
\chi_x\circ \alpha_q^n=\chi_{q^n x}
\qquad(n\ge 0).
\]
因而对任意 \(x,y\in G_S\)，利用字符正交性有
\[
\int_{\Sigma_S}(\chi_x\circ \alpha_q^n)\,\overline{\chi_y}\,d\mu_S
=
\int_{\Sigma_S}\chi_{q^n x-y}\,d\mu_S
=
\begin{cases}
1,& q^n x=y,\\
0,& q^n x\neq y.
\end{cases}
\]

现设 \(q\neq \pm1\)。若存在 \(x,y\neq 0\) 与两个不同整数 \(m>n\) 使
\[
q^m x=y=q^n x,
\]
则
\[
(q^{m-n}-1)x=0.
\]
由于 \(G_S=\ZZ[S^{-1}]\subset\QQ\) 无挠，必有 \(q^{m-n}=1\)，与 \(q\neq \pm1\) 矛盾。故对任意非零 \(x,y\in G_S\)，方程 \(q^n x=y\) 至多成立一次；于是上式对所有充分大的 \(n\) 恒为 \(0\)。

由三角多项式在 \(L^2(\Sigma_S,\mu_S)\) 中稠密，便得到：对任意零均值函数 \(f,g\in L^2(\Sigma_S,\mu_S)\)，有
\[
\langle f\circ \alpha_q^n,g\rangle\longrightarrow 0.
\]
故 \(\alpha_q\) mixing，从而也遍历。

反之，若 \(q=1\)，则 \(\alpha_q=\mathrm{id}\)，显然不遍历。若 \(q=-1\)，取任意非零 \(x\in G_S\)。则
\[
\chi_x\circ \alpha_{-1}^{2n}=\chi_x
\qquad(n\ge 0),
\]
故 \(\alpha_{-1}\) 不是 mixing。并且
\[
\Re\chi_x\circ \alpha_{-1}
=
\Re\chi_{-x}
=
\Re\chi_x,
\]
而 \(\Re\chi_x\) 非常值，故 \(\alpha_{-1}\) 亦不遍历。综上得到所示等价关系。证毕。
\end{proof}

\begin{corollary}[同一指数主率下存在无穷多 pairwise non-isomorphic 的 mixing 局部化系统]\label{cor:xi-time-part69-same-exponential-rate-infinite-mixing-family}
\leanverified{paper\_xi\_time\_part69\_same\_exponential\_rate\_infinite\_mixing\_family}
固定任意既约有理数
\[
q=\frac{a}{b}>1,
\qquad
a,b\in\NN,
\qquad
\gcd(a,b)=1.
\]
对任意有限素数集
\[
S\supseteq \operatorname{Supp}(ab),
\]
记
\[
\Sigma_S:=\widehat{\ZZ[S^{-1}]},
\qquad
\alpha_q:\Sigma_S\to\Sigma_S
\]
为与离散端乘法 \(x\mapsto qx\) 对偶的连续群自同构。则：
\begin{enumerate}
  \item 对所有这类 \(S\)，都有
  \[
  \lim_{m\to\infty}\frac1m\log \#\operatorname{Fix}(\alpha_q^m)=\log a;
  \]
  \item 若 \(S\neq T\)，则
  \[
  \Sigma_S\not\cong \Sigma_T;
  \]
  \item 因而对每个固定的 \(q>1\)，存在无穷多个 pairwise non-isomorphic 的 compact connected mixing 系统
  \[
  (\Sigma_S,\alpha_q),
  \]
  它们具有完全相同的周期指数主率 \(\log a\)。
\end{enumerate}
\end{corollary}
\begin{proof}
第 \(1\) 条是定理 \ref{thm:xi-time-part69-unit-automorphism-polynomial-primeblindness} 的直接特例，因为此时
\[
\max\{|a|,b\}=a.
\]
第 \(2\) 条由定理 \ref{thm:xi-time-part56-localized-pontryagin-isomorphism-rigidity} 立即得到。

对第 \(3\) 条，只需注意：包含 \(\operatorname{Supp}(ab)\) 的有限素数集 \(S\) 有无穷多个；由定理 \ref{thm:xi-time-part69-unit-automorphism-ergodic-mixing-collapse}，对每个这样的 \(S\)，因 \(q>1\) 故 \(\alpha_q\) mixing；再结合前两条，便得到一列两两不同构而指数主率完全相同的局部化动力系统。证毕。
\end{proof}

\endinput

\paragraph{有限角色审计的单圆塌缩与 prime-register 尾部失明}
在局部化 solenoid 的连续自同构、周期点指数与 prime-register 的多项式级盲性已经闭合之后，有限角色族的审计能力还可进一步压缩为一条双重塌缩链：在连通方向上，任意有限角色族都退化为单一 primitive 圆相位；在全不连通方向上，它们对 prime-register 核只能读取有限深度截断，因此有限 trigonometric 审计代数同时失去高维相位自由度与无限 \(p\)-进尾部。

\begin{theorem}[局部化 solenoid 上任意有限角色族都塌缩为单一 primitive 相位]\label{thm:xi-time-part69-finite-character-family-single-primitive-phase}
\leanverified{paper\_xi\_time\_part69\_finite\_character\_family\_single\_primitive\_phase}
设 \(S\subset\mathbb P\) 为有限素数集，并记
\[
G_S:=\ZZ[S^{-1}],
\qquad
\Sigma_S:=\widehat{G_S}.
\]
任取有限族连续角色
\[
\chi_1,\dots,\chi_r\in \widehat{\Sigma_S}\cong G_S.
\]
则存在一个角色 \(\chi_\ast\in\widehat{\Sigma_S}\) 与整数
\[
n_1,\dots,n_r\in\ZZ
\]
使得
\[
\chi_j=\chi_\ast^{\,n_j}
\qquad(1\le j\le r).
\]
若 \(\chi_1,\dots,\chi_r\) 不全为平凡角色，则 \(\chi_\ast\) 可取为 primitive 角色，并且在此规范下唯一到反演；相应地，整数族 \((n_1,\dots,n_r)\) 自动满足
\[
\gcd(n_1,\dots,n_r)=1.
\]
\end{theorem}
\begin{proof}
在 Pontryagin 对偶同构
\(
\widehat{\Sigma_S}\cong G_S
\)
下，令 \(\chi_j\) 对应于 \(q_j\in G_S\)。记
\[
H:=\langle q_1,\dots,q_r\rangle\subseteq G_S\subseteq \QQ.
\]
取正整数 \(D\) 使得
\[
Dq_j\in\ZZ
\qquad(1\le j\le r).
\]
则 \(DH\) 是 \(\ZZ\) 的有限生成子群，故必为循环群。于是存在唯一整数 \(d\ge 0\) 使
\[
DH=d\ZZ.
\]

若 \(d=0\)，则 \(H=\{0\}\)，亦即全部 \(q_j=0\)，从而所有 \(\chi_j\) 都平凡。此时取 \(\chi_\ast=1\) 与 \(n_j=0\) 即可。

若 \(d>0\)，则
\[
H=\frac{d}{D}\ZZ.
\]
记
\[
q_\ast:=\frac{d}{D}\in G_S.
\]
则 \(H=\ZZ q_\ast\)，故存在整数 \(n_j\) 使
\[
q_j=n_j q_\ast
\qquad(1\le j\le r).
\]
令 \(\chi_\ast\) 为与 \(q_\ast\) 对应的角色，则
\[
\chi_j=\chi_{q_j}=\chi_{q_\ast}^{\,n_j}=\chi_\ast^{\,n_j}.
\]
若 \(q_\ast'\) 也是 \(H\) 的生成元，则 \(q_\ast'=\pm q_\ast\)，故对应角色只差反演 \(\chi_\ast^{-1}\)。最后，由
\[
H=\langle q_1,\dots,q_r\rangle=\langle n_1 q_\ast,\dots,n_r q_\ast\rangle
\]
可知 \(\langle n_1,\dots,n_r\rangle=\ZZ\)，从而 \(\gcd(n_1,\dots,n_r)=1\)。证毕。
\end{proof}

\begin{corollary}[有限相位审计的全部 Borel 结构都等价于单圆相位审计]\label{cor:xi-time-part69-finite-phase-audit-borel-collapse-single-circle}
沿用定理 \ref{thm:xi-time-part69-finite-character-family-single-primitive-phase} 的记号，并定义联合相位映射
\[
\boldsymbol{\chi}:\Sigma_S\to \TT^r,
\qquad
x\longmapsto \bigl(\chi_1(x),\dots,\chi_r(x)\bigr).
\]
则存在一个角色 \(\chi_\ast:\Sigma_S\to\TT\) 与整数指数单项式映射
\[
M:\TT\to \TT^r,
\qquad
M(z):=(z^{n_1},\dots,z^{n_r}),
\]
使得
\[
\boldsymbol{\chi}=M\circ \chi_\ast.
\]
从而对任意 Borel 集 \(B\subseteq \TT^r\)，都存在 Borel 集 \(B_\ast\subseteq \TT\) 使
\[
\boldsymbol{\chi}^{-1}(B)=\chi_\ast^{-1}(B_\ast).
\]
并且有限角色族所生成的 \(\sigma\)-代数恰满足
\[
\sigma(\chi_1,\dots,\chi_r)=\sigma(\chi_\ast).
\]
\end{corollary}
\begin{proof}
由定理 \ref{thm:xi-time-part69-finite-character-family-single-primitive-phase}，存在 \(\chi_\ast\) 与整数 \(n_1,\dots,n_r\) 使
\[
\chi_j=\chi_\ast^{\,n_j}.
\]
故对任意 \(x\in \Sigma_S\)，有
\[
\boldsymbol{\chi}(x)
=
\bigl(\chi_\ast(x)^{n_1},\dots,\chi_\ast(x)^{n_r}\bigr)
=
M(\chi_\ast(x)),
\]
即 \(\boldsymbol{\chi}=M\circ \chi_\ast\)。

任取 Borel 集 \(B\subseteq \TT^r\)，令
\[
B_\ast:=M^{-1}(B)\subseteq \TT.
\]
由于 \(M\) 连续，\(B_\ast\) 仍为 Borel 集，并且
\[
\boldsymbol{\chi}^{-1}(B)
=
\chi_\ast^{-1}(M^{-1}(B))
=
\chi_\ast^{-1}(B_\ast).
\]
这给出
\[
\sigma(\chi_1,\dots,\chi_r)\subseteq \sigma(\chi_\ast).
\]

若全部 \(\chi_j\) 平凡，则两侧 \(\sigma\)-代数都只含空集与全集，结论已得。下设并非全部平凡。由定理 \ref{thm:xi-time-part69-finite-character-family-single-primitive-phase}，
\[
\gcd(n_1,\dots,n_r)=1.
\]
故可取整数 \(a_1,\dots,a_r\) 使
\[
\sum_{j=1}^{r}a_j n_j=1.
\]
定义连续映射
\[
N:\TT^r\to \TT,
\qquad
N(z_1,\dots,z_r):=\prod_{j=1}^{r}z_j^{a_j}.
\]
则对任意 \(x\in \Sigma_S\)，有
\[
N(\boldsymbol{\chi}(x))
=
\prod_{j=1}^{r}\chi_j(x)^{a_j}
=
\chi_\ast(x)^{\sum_j a_j n_j}
=
\chi_\ast(x).
\]
于是
\[
\chi_\ast=N\circ \boldsymbol{\chi}.
\]
对任意 Borel 集 \(C\subseteq \TT\)，有
\[
\chi_\ast^{-1}(C)=\boldsymbol{\chi}^{-1}(N^{-1}(C)),
\]
从而
\[
\sigma(\chi_\ast)\subseteq \sigma(\chi_1,\dots,\chi_r).
\]
综上得到两侧 \(\sigma\)-代数相等。证毕。
\end{proof}

\begin{theorem}[有限角色族对 prime-register 核只能看到有限层截断]\label{thm:xi-time-part69-finite-character-prime-ledger-finite-depth}
\leanverified{paper\_xi\_time\_part69\_finite\_character\_prime\_ledger\_finite\_depth}
沿用定理 \ref{thm:cdim-star-z1s-dual-extension} 的短正合列
\[
0\longrightarrow K_S:=\prod_{p\in S}\ZZ_p
\longrightarrow \Sigma_S
\longrightarrow \TT
\longrightarrow 0.
\]
给定任意有限角色族 \(\chi_1,\dots,\chi_r\in \widehat{\Sigma_S}\cong G_S\)，令 \(q_j\in G_S\) 为其对应元素，并对每个 \(p\in S\) 定义
\[
m_{j,p}:=\max\{0,-v_p(q_j)\}.
\]
当 \(q_j\neq 0\) 时，等价地，\(q_j\) 可唯一写成
\[
q_j=a_j\prod_{p\in S}p^{-m_{j,p}},
\qquad
a_j\in\ZZ,
\qquad
v_p(a_j)=0\ \ (p\in S).
\]
若 \(q_j=0\)，则取全部 \(m_{j,p}=0\)。
再令
\[
M_p:=\max_{1\le j\le r}m_{j,p},
\qquad
U(\chi_1,\dots,\chi_r):=\prod_{p\in S}p^{M_p}\ZZ_p\subseteq K_S.
\]
则：
\begin{enumerate}
  \item 每个限制角色 \(\chi_j|_{K_S}\) 都在 \(U(\chi_1,\dots,\chi_r)\) 的陪集上恒定；
  \item 若记
  \[
  T_{\boldsymbol{\chi}}:=\{p\in S:\ M_p\ge 1\},
  \]
  则联合映射
  \[
  K_S\longrightarrow \TT^r,
  \qquad
  x\longmapsto \bigl(\chi_1(x),\dots,\chi_r(x)\bigr)
  \]
  必经由有限商
  \[
  K_S/U(\chi_1,\dots,\chi_r)
  \cong
  \prod_{p\in T_{\boldsymbol{\chi}}}\ZZ/p^{M_p}\ZZ
  \]
  因子化；
  \item 因而任何有限角色审计都严格失明于 \(K_S\) 的无限 \(p\)-进尾部。
\end{enumerate}
\end{theorem}
\begin{proof}
由定理 \ref{thm:cdim-star-z1s-dual-extension}，有互为 Pontryagin 对偶的短正合列
\[
0\longrightarrow \ZZ
\longrightarrow G_S
\longrightarrow \bigoplus_{p\in S}C_{p^\infty}
\longrightarrow 0
\]
与
\[
0\longrightarrow K_S
\longrightarrow \Sigma_S
\longrightarrow \TT
\longrightarrow 0.
\]
因此，\(\chi_j\) 在 \(K_S\) 上的限制正对应于 \(q_j\) 在商群
\[
G_S/\ZZ\cong \bigoplus_{p\in S}C_{p^\infty}
\]
中的像。

固定 \(j\) 与 \(p\in S\)。由 \(m_{j,p}=\max\{0,-v_p(q_j)\}\) 的定义，\(q_j\) 在 \(G_S/\ZZ\) 中的 \(p\)-主分量落在 \(C_{p^\infty}\) 的唯一 \(p^{m_{j,p}}\)-阶循环子群内。对偶化以后，这意味着 \(\chi_j|_{K_S}\) 的第 \(p\)-坐标分量只能通过有限商
\[
\ZZ_p\twoheadrightarrow \ZZ_p/p^{m_{j,p}}\ZZ_p\cong \ZZ/p^{m_{j,p}}\ZZ
\]
被读取；等价地，
\[
p^{m_{j,p}}\ZZ_p
\subseteq
\ker\bigl(\chi_j|_{\ZZ_p}\bigr).
\]
于是定义
\[
U_j:=\prod_{p\in S}p^{m_{j,p}}\ZZ_p\subseteq K_S,
\]
便有
\[
U_j\subseteq \ker(\chi_j|_{K_S}).
\]

对全部 \(j\) 取交，得到
\[
\bigcap_{j=1}^{r}U_j
=
\prod_{p\in S}p^{M_p}\ZZ_p
=
U(\chi_1,\dots,\chi_r).
\]
故每个 \(\chi_j|_{K_S}\) 都在 \(U(\chi_1,\dots,\chi_r)\) 的陪集上恒定，第 \(1\) 条成立。

由于 \(U(\chi_1,\dots,\chi_r)\) 是 \(K_S\) 的开子群，联合映射必经由商群 \(K_S/U(\chi_1,\dots,\chi_r)\) 因子化。又由各坐标逐素数分解，
\[
K_S/U(\chi_1,\dots,\chi_r)
\cong
\prod_{p\in T_{\boldsymbol{\chi}}}\ZZ/p^{M_p}\ZZ.
\]
这给出第 \(2\) 条。最后，所有更深层的 \(p\)-进尾部都落在 \(U(\chi_1,\dots,\chi_r)\) 内，因而不会改变任何有限角色的取值，第 \(3\) 条随即成立。证毕。
\end{proof}

\begin{theorem}[有限角色生成的 trigonometric 审计代数在一致闭包后恰为单圆代数]\label{thm:xi-time-part69-finite-trigonometric-audit-algebra-single-circle-finite-tail}
设
\[
\mathcal A(\chi_1,\dots,\chi_r)\subset C(\Sigma_S)
\]
为由有限角色族 \(\chi_1,\dots,\chi_r\) 及其复共轭生成的含幺 \(^\ast\)-代数。令 \(\chi_\ast\) 为定理 \ref{thm:xi-time-part69-finite-character-family-single-primitive-phase} 给出的生成角色，并令
\[
U:=U(\chi_1,\dots,\chi_r)\subseteq K_S
\]
为定理 \ref{thm:xi-time-part69-finite-character-prime-ledger-finite-depth} 的开子群。则成立：
\begin{enumerate}
  \item 代数恒等式
  \[
  \mathcal A(\chi_1,\dots,\chi_r)
  =
  \chi_\ast^{*}\bigl(\CC[z,z^{-1}]\bigr);
  \]
  \item 若 \(\chi_1,\dots,\chi_r\) 不全平凡，则
  \[
  \overline{\mathcal A(\chi_1,\dots,\chi_r)}^{\|\cdot\|_\infty}
  =
  \chi_\ast^{*}C(\TT)
  \cong
  C(\TT);
  \]
  若全部平凡，则
  \[
  \overline{\mathcal A(\chi_1,\dots,\chi_r)}^{\|\cdot\|_\infty}
  =
  \CC\mathbf 1;
  \]
  \item 由 \(\mathcal A(\chi_1,\dots,\chi_r)\) 生成的任意有限 Borel 分割的每个非空原子都包含某个 \(U\)-陪集；特别地，它不可能把 \(K_S\) 的点分离到有限精度以下。
\end{enumerate}
\end{theorem}
\begin{proof}
由定理 \ref{thm:xi-time-part69-finite-character-family-single-primitive-phase}，存在整数 \(n_1,\dots,n_r\) 使
\[
\chi_j=\chi_\ast^{\,n_j}.
\]
因此 \(\mathcal A(\chi_1,\dots,\chi_r)\) 的每个元素都是 \(\chi_\ast\) 的 Laurent 多项式，故有包含
\[
\mathcal A(\chi_1,\dots,\chi_r)
\subseteq
\chi_\ast^{*}\bigl(\CC[z,z^{-1}]\bigr).
\]

若全部 \(\chi_j\) 平凡，则两边都等于常数代数，结论立即成立。下设并非全部平凡。由定理 \ref{thm:xi-time-part69-finite-character-family-single-primitive-phase}，可取 \(\chi_\ast\) primitive，且
\[
\gcd(n_1,\dots,n_r)=1.
\]
故可取整数 \(a_1,\dots,a_r\) 使
\[
\sum_{j=1}^{r}a_j n_j=1.
\]
于是
\[
\chi_\ast=\prod_{j=1}^{r}\chi_j^{a_j}.
\]
由于负指数可由复共轭实现，故 \(\chi_\ast\in \mathcal A(\chi_1,\dots,\chi_r)\)。从而
\[
\chi_\ast^{*}\bigl(\CC[z,z^{-1}]\bigr)
\subseteq
\mathcal A(\chi_1,\dots,\chi_r),
\]
第 \(1\) 条得证。

对第 \(2\) 条，只需证明非平凡情形下 \(\chi_\ast\) 满射到 \(\TT\)。由于 \(G_S\) 无挠，\(\Sigma_S=\widehat{G_S}\) 是紧连通交换群；故 \(\chi_\ast(\Sigma_S)\) 是 \(\TT\) 的闭连通子群。又 \(\chi_\ast\) 非平凡，故其像不能为 \(\{1\}\)，从而
\[
\chi_\ast(\Sigma_S)=\TT.
\]
于是 \(\chi_\ast^{*}:C(\TT)\to C(\Sigma_S)\) 是等距单射。由 Stone--Weierstrass 定理，
\[
\overline{\CC[z,z^{-1}]}^{\|\cdot\|_\infty}=C(\TT),
\]
拉回后即得
\[
\overline{\mathcal A(\chi_1,\dots,\chi_r)}^{\|\cdot\|_\infty}
=
\chi_\ast^{*}C(\TT)
\cong
C(\TT).
\]

最后，由定理 \ref{thm:xi-time-part69-finite-character-prime-ledger-finite-depth}，每个生成元 \(\chi_j\) 都在 \(U\)-陪集上恒定。故 \(\mathcal A(\chi_1,\dots,\chi_r)\) 的每个元素，以及其一致闭包中的每个元素，也都在 \(U\)-陪集上恒定。于是由这些函数所生成的全部 Borel 集都是 \(U\)-饱和集，即都是 \(U\)-陪集的并。因而任意由该代数决定的有限 Borel 分割，其每个非空原子也必为 \(U\)-陪集的并；特别地，取其中任一点 \(x\)，整个陪集 \(xU\) 都包含在该原子中。故该分割不可能把 \(K_S\) 的点分离到比 \(U\) 更细的精度。证毕。
\end{proof}

\endinput

\paragraph{有限连续可见性、有限指数/有限商谱与连续自同态环的闭合压缩}
在局部化数系的对偶短正合列、有限商分类、连续自同态与有限角色审计已经分别获得闭式之后，还可把连续有限可见性、有限指数格、有限商轮廓与连续动力学几条线路进一步压缩为下列若干条直接后果。

\begin{theorem}[局部化数系对偶在有限连续读出下的完全塌缩]\label{thm:xi-time-part69c-localized-dual-finite-visible-collapse}
设 \(S\subset\mathbb P\) 为有限素数集，并记
\[
G_S:=\ZZ[S^{-1}],
\qquad
\Sigma_S:=\widehat{G_S}.
\]
则 \(\Sigma_S\) 是连通紧交换群。因而对任意有限离散集 \(A\) 与任意连续映射
\[
F:\Sigma_S\longrightarrow A,
\]
映射 \(F\) 必为常值。
\end{theorem}
\begin{proof}
由定理 \ref{thm:cdim-star-z1s-dual-extension}，\(\Sigma_S\) 为紧交换群，并落在短正合列
\[
0\longrightarrow \prod_{p\in S}\ZZ_p
\longrightarrow \Sigma_S
\longrightarrow \TT
\longrightarrow 0.
\]
另一方面，
\[
G_S=\ZZ[S^{-1}]\subset \QQ
\]
作为离散加法群无挠。按 Pontryagin 对偶的标准连通性判据，\(\Sigma_S\) 必连通。

现令 \(A\) 为有限离散集，\(F:\Sigma_S\to A\) 连续。则 \(F(\Sigma_S)\) 是 \(A\) 的连通子空间；而有限离散空间的连通子空间只能是单点，故 \(F(\Sigma_S)\) 仅含一个元素，从而 \(F\) 为常值。证毕。
\end{proof}

\begin{theorem}[局部化数系有限商阶谱的完全闭式]\label{thm:xi-time-part69c-localized-finite-quotient-order-spectrum}
\leanverified{paper\_xi\_time\_part69c\_localized\_finite\_quotient\_order\_spectrum}
沿用定理 \ref{thm:xi-time-part69c-localized-dual-finite-visible-collapse} 的记号，并定义
\[
\mathcal Q(G_S)
:=
\left\{
n\ge 1:\ \exists\ \text{满同态 }G_S\twoheadrightarrow \ZZ/n\ZZ
\right\}.
\]
则有严格等式
\[
\mathcal Q(G_S)
=
\left\{
n\ge 1:\ \gcd\!\left(n,\prod_{p\in S}p\right)=1
\right\}.
\]
等价地，\(G_S\) 的全部有限商恰为循环群
\[
\ZZ/n\ZZ
\qquad
\bigl(\gcd(n,\prod_{p\in S}p)=1\bigr).
\]
\end{theorem}
\begin{proof}
由定理 \ref{thm:xi-localized-finite-quotient-phase-layer-classification} 第 \(1\) 条，对任意 \(n\ge 1\) 都有
\[
G_S/nG_S\cong \ZZ/n_{S^\perp}\ZZ.
\]
若
\[
\gcd\!\left(n,\prod_{p\in S}p\right)=1,
\]
则 \(n_{S^\perp}=n\)，故自然商映射给出满同态
\[
G_S\twoheadrightarrow G_S/nG_S\cong \ZZ/n\ZZ,
\]
于是 \(n\in \mathcal Q(G_S)\)。

反之，若 \(n\in \mathcal Q(G_S)\)，则存在满同态
\[
G_S\twoheadrightarrow \ZZ/n\ZZ.
\]
这说明 \(\ZZ/n\ZZ\) 是 \(G_S\) 的有限商。再由定理
\ref{thm:xi-localized-finite-quotient-phase-layer-classification}
第 \(2\) 条，\(G_S\) 的每个有限商都必为阶数与 \(S\) 中所有素数互素的循环群。故 \(n\) 的全部素因子都不属于 \(S\)，亦即
\[
\gcd\!\left(n,\prod_{p\in S}p\right)=1.
\]
于是所示集合恒等式成立；最后一句只是同一定理的重述。证毕。
\end{proof}

\begin{theorem}[局部化数系的有限指数子群完全分类]\label{thm:xi-time-part69c-localized-finite-index-subgroup-classification}
沿用定理 \ref{thm:xi-time-part69c-localized-finite-quotient-order-spectrum} 的记号。则 \(G_S\) 的任意有限指数子群 \(H\le G_S\) 都唯一形如
\[
H=nG_S
\]
其中
\[
n\ge 1,
\qquad
\gcd\!\left(n,\prod_{p\in S}p\right)=1.
\]
并且其指数满足精确公式
\[
[G_S:nG_S]=n.
\]
\leanverified{paper\_xi\_time\_part69c\_localized\_finite\_index\_subgroup\_classification}
\end{theorem}
\begin{proof}
先设
\[
\gcd\!\left(n,\prod_{p\in S}p\right)=1.
\]
定义模 \(n\) 约化映射
\[
\rho_n:G_S\to \ZZ/n\ZZ,
\qquad
\rho_n\!\left(\frac{a}{\prod_{p\in S}p^{k_p}}\right)
:=
a\cdot\left(\prod_{p\in S}p^{k_p}\right)^{-1}\pmod n.
\]
由于每个 \(p\in S\) 在 \(\ZZ/n\ZZ\) 中可逆，\(\rho_n\) 良定义。它显然满射，因为整数子群 \(\ZZ\subset G_S\) 已满射到 \(\ZZ/n\ZZ\)。

现证
\[
\ker(\rho_n)=nG_S.
\]
包含 \(nG_S\subseteq \ker(\rho_n)\) 直接成立。反向地，若
\[
x=\frac{a}{b}\in G_S,
\qquad
b=\prod_{p\in S}p^{k_p},
\]
且 \(\rho_n(x)=0\)，则
\[
a\,b^{-1}\equiv 0\pmod n.
\]
由于 \(\gcd(b,n)=1\)，\(b^{-1}\) 在 \(\ZZ/n\ZZ\) 中可逆，故 \(a\equiv 0\pmod n\)。写 \(a=nc\)，便有
\[
x=\frac{nc}{b}=n\cdot \frac{c}{b}\in nG_S.
\]
因此
\[
G_S/nG_S\cong \ZZ/n\ZZ,
\]
从而
\[
[G_S:nG_S]=n.
\]

反之，设 \(H\le G_S\) 为有限指数子群，并记
\[
A:=G_S/H.
\]
则 \(A\) 为有限商。由定理
\ref{thm:xi-time-part69c-localized-finite-quotient-order-spectrum}，存在唯一满足
\[
\gcd\!\left(n,\prod_{p\in S}p\right)=1
\]
的正整数 \(n\)，使得
\[
A\cong \ZZ/n\ZZ.
\]
因而 \([G_S:H]=n\)，并且 \(nA=0\)。这说明对每个 \(x\in G_S\) 都有
\[
nx\in H,
\]
故
\[
nG_S\subseteq H.
\]
而第一部分已证明
\[
[G_S:nG_S]=n=[G_S:H].
\]
既然 \(nG_S\subseteq H\) 且两者指数相同，遂得
\[
H=nG_S.
\]
唯一性则由指数唯一确定。证毕。
\end{proof}

\begin{corollary}[有限指数子群格与 \(S\)-自由整除格的反同构及其子群 \texorpdfstring{$\zeta$}{zeta} 函数]\label{cor:xi-time-part69c-localized-finite-index-lattice-zeta}
沿用定理 \ref{thm:xi-time-part69c-localized-finite-index-subgroup-classification} 的记号。则映射
\[
n\longmapsto nG_S
\qquad
\left(
\gcd\!\left(n,\prod_{p\in S}p\right)=1
\right)
\]
把 \(S\)-自由正整数的整除偏序反同构到 \(G_S\) 的有限指数子群格。更具体地，对任意与 \(S\) 互素的正整数 \(m,n\)，有
\[
nG_S\subseteq mG_S
\quad\Longleftrightarrow\quad
m\mid n,
\]
并且
\[
nG_S\cap mG_S=\operatorname{lcm}(n,m)\,G_S,
\]
\[
nG_S+mG_S=\gcd(n,m)\,G_S.
\]
因此 \(G_S\) 的子群 \(\zeta\) 函数满足完全闭式
\[
\zeta_{G_S}^{\le}(s)
:=
\sum_{H\le_f G_S}[G_S:H]^{-s}
=
\sum_{\gcd(n,\prod_{p\in S}p)=1}n^{-s}
=
\zeta(s)\prod_{p\in S}(1-p^{-s}).
\]
\end{corollary}
\begin{proof}
先证包含判据。若 \(m\mid n\)，写 \(n=mk\)，则
\[
nG_S=m(kG_S)\subseteq mG_S,
\]
故 \(nG_S\subseteq mG_S\)。

反之，若 \(nG_S\subseteq mG_S\)，则特别有 \(n\in mG_S\)，故存在 \(q\in G_S\) 使
\[
n=mq.
\]
于是
\[
q=\frac{n}{m}\in G_S.
\]
但 \(n,m\) 都与 \(S\) 互素，故既约分数 \(n/m\) 的分母不含任何 \(S\)-中素数。要使 \(n/m\in G_S=\ZZ[S^{-1}]\)，其分母只能为 \(1\)。因此 \(n/m\in \ZZ\)，即 \(m\mid n\)。

现设
\[
K:=nG_S\cap mG_S.
\]
由定理 \ref{thm:xi-time-part69c-localized-finite-index-subgroup-classification}，存在唯一 \(S\)-自由正整数 \(k\) 使
\[
K=kG_S.
\]
由 \(K\subseteq nG_S\) 与 \(K\subseteq mG_S\)，按上面的包含判据得
\[
n\mid k,
\qquad
m\mid k,
\]
故 \(\operatorname{lcm}(n,m)\mid k\)。另一方面，
\[
\operatorname{lcm}(n,m)\,G_S\subseteq nG_S\cap mG_S=K,
\]
再次应用包含判据得
\[
k\mid \operatorname{lcm}(n,m).
\]
故 \(k=\operatorname{lcm}(n,m)\)，从而
\[
nG_S\cap mG_S=\operatorname{lcm}(n,m)\,G_S.
\]

再设
\[
L:=nG_S+mG_S.
\]
同理，\(L=\ell G_S\) 对某个唯一 \(S\)-自由正整数 \(\ell\) 成立。由
\[
nG_S\subseteq L,
\qquad
mG_S\subseteq L,
\]
可知
\[
\ell\mid n,
\qquad
\ell\mid m,
\]
故 \(\ell\mid \gcd(n,m)\)。另一方面，取整数 \(u,v\) 满足
\[
un+vm=\gcd(n,m).
\]
则对任意 \(x\in G_S\)，有
\[
\gcd(n,m)\,x=u(nx)+v(mx)\in nG_S+mG_S=L.
\]
因此
\[
\gcd(n,m)\,G_S\subseteq L=\ell G_S,
\]
故由包含判据得
\[
\ell\mid \gcd(n,m).
\]
另一方面，由 \(n,m\) 都是 \(\gcd(n,m)\) 的整数倍，显然
\[
nG_S\subseteq \gcd(n,m)\,G_S,
\qquad
mG_S\subseteq \gcd(n,m)\,G_S.
\]
从而
\[
L=nG_S+mG_S\subseteq \gcd(n,m)\,G_S,
\]
再次应用包含判据得到
\[
\gcd(n,m)\mid \ell.
\]
于是 \(\ell=\gcd(n,m)\)，从而
\[
nG_S+mG_S=\gcd(n,m)\,G_S.
\]

最后，由定理 \ref{thm:xi-time-part69c-localized-finite-index-subgroup-classification}，全部有限指数子群恰由
\[
H=nG_S
\qquad
\left(
\gcd\!\left(n,\prod_{p\in S}p\right)=1
\right)
\]
枚举，且其指数正为 \(n\)。故
\[
\zeta_{G_S}^{\le}(s)
=
\sum_{\gcd(n,\prod_{p\in S}p)=1}n^{-s}.
\]
按 Euler 乘积展开即得
\[
\sum_{\gcd(n,\prod_{p\in S}p)=1}n^{-s}
=
\prod_{p\notin S}(1-p^{-s})^{-1}
=
\zeta(s)\prod_{p\in S}(1-p^{-s}).
\]
证毕。
\end{proof}

\begin{corollary}[有限商轮廓对素数账本的完全恢复]\label{cor:xi-time-part69c-localized-finite-quotient-profile-recovers-primes}
对任意有限素数集 \(S,T\subset\mathbb P\)，有
\[
\mathcal Q(G_S)=\mathcal Q(G_T)
\quad\Longleftrightarrow\quad
S=T.
\]
等价地，对任意素数 \(p\) 都有
\[
p\in S
\quad\Longleftrightarrow\quad
\forall n\in \mathcal Q(G_S),\ p\nmid n.
\]
\end{corollary}
\begin{proof}
由定理 \ref{thm:xi-time-part69c-localized-finite-quotient-order-spectrum}，
\[
\mathcal Q(G_S)
=
\left\{
n\ge 1:\ \forall p\in S,\ p\nmid n
\right\}.
\]
因此 \(\mathcal Q(G_S)\) 恰记录了哪些素数被完全排除在全部有限商阶谱之外。若
\[
\mathcal Q(G_S)=\mathcal Q(G_T),
\]
则被排除的素数集合一致，故 \(S=T\)。反向蕴含显然。

第二条只是同一刻画按单个素数逐点展开：若 \(p\in S\)，则由上式 \(p\) 不整除任意 \(n\in\mathcal Q(G_S)\)；反之若 \(p\notin S\)，则 \(n=p\) 已属于 \(\mathcal Q(G_S)\)。证毕。
\end{proof}

\begin{theorem}[局部化数系的连续自同态环完全由局部化环穷尽]\label{thm:xi-time-part69c-localized-dual-endomorphism-ring-exhaustion}
沿用定理 \ref{thm:xi-time-part69c-localized-dual-finite-visible-collapse} 的记号。则加法群 \(G_S\) 的端同态环满足
\[
\operatorname{End}(G_S)\cong \ZZ[S^{-1}],
\]
其中同构由
\[
q\longmapsto m_q,
\qquad
m_q(x):=qx
\]
给出。因而其对偶紧群的连续端同态环满足
\[
\operatorname{End}_{\mathrm{cts}}(\Sigma_S)\cong \ZZ[S^{-1}].
\]
进一步，连续自同构群满足
\[
\operatorname{Aut}_{\mathrm{cts}}(\Sigma_S)
\cong
\operatorname{Aut}(G_S)
\cong
\left\{
\pm\prod_{p\in S}p^{n_p}: n_p\in\ZZ
\right\}.
\]
\end{theorem}
\begin{proof}
定理 \ref{thm:xi-localized-endomorphism-ring-solenoid-indecomposable} 的前两条已经给出
\[
\operatorname{End}(G_S)\cong G_S=\ZZ[S^{-1}],
\]
并且该同构正由乘法映射 \(q\mapsto m_q\) 实现。Pontryagin 对偶把紧交换群上的连续端同态环识别为离散端端同态环的反环，故
\[
\operatorname{End}_{\mathrm{cts}}(\Sigma_S)\cong \operatorname{End}(G_S)^{\mathrm{op}}.
\]
由于 \(\ZZ[S^{-1}]\) 为交换环，其反环与自身一致，于是
\[
\operatorname{End}_{\mathrm{cts}}(\Sigma_S)\cong \ZZ[S^{-1}].
\]

最后，由定理 \ref{thm:xi-time-part69-localized-unit-automorphism-group-classification}，
\[
\operatorname{Aut}_{\mathrm{cts}}(\Sigma_S)
\cong
G_S^\times
=
\left\{
\pm\prod_{p\in S}p^{n_p}: n_p\in\ZZ
\right\}.
\]
而离散端自同构正对应于 \(\ZZ[S^{-1}]\) 的可逆乘法元，故同一公式也描述 \(\operatorname{Aut}(G_S)\)。证毕。
\end{proof}

\noindent
由此，局部化数系的对偶对象在连续有限值读出层面完全塌缩为单点；而在有限指数与有限商层面，其全部可见结构又分别被精确压缩为 \(S\)-自由整数的整除半格、对应的子群 \(\zeta\) 函数 Euler 因子以及有限商阶谱。进一步，在连续自同态动力学层面，全部端同态与自同构又完全坍缩为局部化环 \(\ZZ[S^{-1}]\) 及其单位群本身。于是，素数账本、有限指数格、有限商轮廓与连续动力学四者被同一局部化数据闭合地锁定。

\input{subsubsec__xi-time-protocol-conclusions-part69d-localized-hom-godeltate-centralizer-periodic}

\endinput

\paragraph{六模拓扑相图、精确暗块、boundary 奇偶盲性与纯鸽巢截断}
在最大非可缩纤维的模 \(6\) 比例极限、审计偶窗最小退化扇区的奇偶同伦相型、window-\(6\) 的 boundary sheet-parity 中心分裂，以及黄金分支在 Fibonacci 分母上的 \(\mathsf T\)/\(D^\ast\) 严格奇偶双相都已固定之后，还可进一步抽取出下列四条跨链后果。它们分别刻画精确暗块与顶端非可缩纤维的非耦合共存、根字典对 boundary 奇偶三格的严格盲性、审计偶窗最弱退化扇区在阿贝尔化中的精确 Fibonacci 份额，以及最小退化 Fibonacci 律一旦试图脱离审计窗口便会遭遇的首个纯计数截断。

\begin{corollary}[精确暗块与拓扑主纤维的非耦合共存]\label{cor:xi-time-part69c-darkblock-topfiber-decoupling}
沿用推论 \ref{cor:xi-binfold-zero-block-collision-gap-decoupling} 与定理 \ref{thm:xi-time-part62db-maxfiber-homotopy-phase-mod6} 的记号，记
\[
D_m:=\max_{x\in X_m}d_m(x),
\qquad
\widetilde D_m:=\max\{d_m(x):x\in X_m,\ |K_x|\not\simeq *\},
\qquad
\bar d_m:=\frac{2^m}{F_{m+2}}.
\]
若
\[
m\equiv 4\pmod 6,
\]
则同时成立
\[
|\mathcal Z_m|\ge F_{(m+2)/2},
\qquad
\widetilde D_m=D_m,
\]
并且
\[
\liminf_{\substack{m\to\infty\\ m\equiv 4\!\!\!\!\pmod 6}}
\frac{D_m}{\bar d_m}
\ge
1+C_\varphi.
\]
因此，半阶指数规模的精确暗块与不发生拓扑折损的顶端非可缩纤维，可以在同一分辨率相上严格共存。
\end{corollary}
\begin{proof}
当 \(m\equiv 4\pmod 6\) 时，\(m+2\equiv 0\pmod 6\)。故由推论 \ref{cor:xi-binfold-zero-block-collision-gap-decoupling}，
\[
|\mathcal Z_m|\ge F_{(m+2)/2}.
\]
同一同余类上，定理 \ref{thm:xi-time-part62db-maxfiber-homotopy-phase-mod6} 直接给出
\[
\widetilde D_m=D_m.
\]
另一方面，结论 \ref{con:xi-time-part9m-cphi-forces-collision-gap-and-peak-bias} 已证明
\[
\liminf_{m\to\infty}\frac{M_m}{\overline d_m}\ge 1+C_\varphi.
\]
将其中记号 \(M_m,\overline d_m\) 分别改写为 \(D_m,\bar d_m\)，再限制到子序列 \(m\equiv 4\pmod 6\)，便得到第三式。证毕。
\end{proof}

\begin{theorem}[window-\texorpdfstring{$6$}{6} 根字典的 boundary 盲性余维恰为 \texorpdfstring{$3$}{3}]\label{thm:xi-time-part69c-window6-root-boundary-blind-codim3}
沿用定理 \ref{thm:xi-window6-cyclic-kernel-bc3-root-datum} 与定理 \ref{thm:xi-window6-center-exact-splitting-bdry-cyclic} 的记号。通过定理 \ref{thm:xi-foldbin-gauge-representation-distribution-law-m6-spectrum} 的直积分解，识别
\[
\bigl(G_6^{\mathrm{bin}}\bigr)^{\mathrm{ab}}
\cong
(\ZZ/2\ZZ)^{X_6}.
\]
令
\[
\rho:
\bigl(G_6^{\mathrm{bin}}\bigr)^{\mathrm{ab}}
\cong
(\ZZ/2\ZZ)^{X_6}
\longrightarrow
(\ZZ/2\ZZ)^{X_6^{\mathrm{cyc}}}
\]
为忘去 \(X_6^{\mathrm{bdry}}\) 三个坐标的自然投影，并记
\[
\iota_{\mathrm{ab}}:
Z\!\bigl(G_6^{\mathrm{bin}}\bigr)
\hookrightarrow
\bigl(G_6^{\mathrm{bin}}\bigr)^{\mathrm{ab}}
\]
为自然交换化嵌入。则
\[
\ker\rho
\cong
(\ZZ/2\ZZ)^3.
\]
更精确地，
\[
\ker\rho
=
\iota_{\mathrm{ab}}(Z_{\partial}),
\qquad
Z_{\partial}
=
\langle \tau_w:\ w\in X_6^{\mathrm{bdry}}\rangle
\subset Z\!\bigl(G_6^{\mathrm{bin}}\bigr).
\]
因此，任何仅经由 \(X_6^{\mathrm{cyc}}\) 上的 \(B_3/C_3\) 根字典读取的可观测量，都必然对这一规范三维 boundary 奇偶子格失明。
\end{theorem}
\begin{proof}
由定理 \ref{thm:xi-foldbin-gauge-representation-distribution-law-m6-spectrum}，
\[
G_6^{\mathrm{bin}}
\cong
\prod_{w\in X_6}\mathfrak S_{d_6(w)}.
\]
由于 \(m=6\) 时全部纤维大小只可能为 \(2,3,4\)，而
\[
\mathfrak S_2^{\mathrm{ab}}\cong
\mathfrak S_3^{\mathrm{ab}}\cong
\mathfrak S_4^{\mathrm{ab}}
\cong
\ZZ/2\ZZ,
\]
故交换化给出规范识别
\[
\bigl(G_6^{\mathrm{bin}}\bigr)^{\mathrm{ab}}
\cong
\prod_{w\in X_6}\mathfrak S_{d_6(w)}^{\mathrm{ab}}
\cong
(\ZZ/2\ZZ)^{X_6}.
\]
再由定理 \ref{thm:xi-window6-cyclic-kernel-bc3-root-datum}，
\[
X_6=X_6^{\mathrm{cyc}}\sqcup X_6^{\mathrm{bdry}},
\qquad
|X_6^{\mathrm{cyc}}|=18,
\qquad
|X_6^{\mathrm{bdry}}|=3.
\]
因此忘去 boundary 三坐标的投影 \(\rho\) 的核恰为
\[
(\ZZ/2\ZZ)^{X_6^{\mathrm{bdry}}}
\cong
(\ZZ/2\ZZ)^3.
\]

另一方面，定理 \ref{thm:xi-window6-center-exact-splitting-bdry-cyclic} 已经把
\[
Z_{\partial}
=
\langle \tau_w:\ w\in X_6^{\mathrm{bdry}}\rangle
\cong
(\ZZ/2\ZZ)^{X_6^{\mathrm{bdry}}}
\]
识别为 \(Z(G_6^{\mathrm{bin}})\) 的规范直和因子。故
\[
\ker\rho=\iota_{\mathrm{ab}}(Z_{\partial})\cong (\ZZ/2\ZZ)^3.
\]
凡仅通过 \(X_6^{\mathrm{cyc}}\) 的根字典读取的可观测量，都必经由 \(\rho\) 因子化，因而在 \(\ker\rho=\iota_{\mathrm{ab}}(Z_{\partial})\) 上必为常值。证毕。
\end{proof}

\begin{corollary}[审计偶窗最弱退化扇区的奇偶份额恰为 \texorpdfstring{$F_m/F_{m+2}$}{F_m/F_{m+2}}]\label{cor:xi-time-part69c-audited-even-minsector-abel-share}
对审计偶窗
\[
m\in\{6,8,10,12\},
\]
记
\[
\mathcal S_{m,*}:=\mathcal S_{m,d_{\min}(m)},
\qquad
\Lambda_m^{\min}:=(\ZZ/2\ZZ)^{\mathcal S_{m,*}}
\subset
\bigl(G_m^{\mathrm{bin}}\bigr)^{\mathrm{ab}},
\]
其中 \(\Lambda_m^{\min}\) 由定理 \ref{thm:xi-audited-even-abel-center-fibonacci-splitting} 的规范直和分裂给出。则
\[
\operatorname{rank}_{\FF_2}\Lambda_m^{\min}=F_m,
\qquad
\operatorname{rank}_{\FF_2}\!\bigl(G_m^{\mathrm{bin}}\bigr)^{\mathrm{ab}}=F_{m+2},
\]
从而
\[
\frac{\operatorname{rank}_{\FF_2}\Lambda_m^{\min}}
{\operatorname{rank}_{\FF_2}\!\bigl(G_m^{\mathrm{bin}}\bigr)^{\mathrm{ab}}}
=
\frac{F_m}{F_{m+2}}.
\]
具体地，
\[
\frac{8}{21},\qquad
\frac{21}{55},\qquad
\frac{55}{144},\qquad
\frac{144}{377}
\]
分别对应 \(m=6,8,10,12\)。
\end{corollary}
\begin{proof}
由定理 \ref{thm:xi-audited-even-abel-center-fibonacci-splitting}，
\[
\bigl(G_m^{\mathrm{bin}}\bigr)^{\mathrm{ab}}
\cong
(\ZZ/2\ZZ)^{\mathcal S_{m,*}}
\oplus
(\ZZ/2\ZZ)^{X_m\setminus\mathcal S_{m,*}},
\]
并且
\[
|\mathcal S_{m,*}|=F_m.
\]
故
\[
\operatorname{rank}_{\FF_2}\Lambda_m^{\min}=F_m.
\]
另一方面，审计偶窗上由定理 \ref{thm:xi-foldbin-audited-even-center-collapse-abel-full-rank}，
\[
\bigl(G_m^{\mathrm{bin}}\bigr)^{\mathrm{ab}}
\cong
(\ZZ/2\ZZ)^{|X_m|}.
\]
再由稳定类型计数律 \(|X_m|=F_{m+2}\)，得到
\[
\operatorname{rank}_{\FF_2}\!\bigl(G_m^{\mathrm{bin}}\bigr)^{\mathrm{ab}}=F_{m+2}.
\]
于是所示比值公式成立，四个具体分数只需代入 \(m=6,8,10,12\) 即得。证毕。
\end{proof}

\begin{theorem}[审计偶窗最小退化 Fibonacci 外推的首个纯计数截断]\label{thm:xi-time-part69c-minsector-fib-extension-pigeonhole-cutoff}
\leanverified{paper\_xi\_time\_part69c\_minsector\_fib\_extension\_pigeonhole\_cutoff}
把审计恒等式
\[
d_{\min}(m)=F_{m/2}
\qquad
\bigl(m\in\{6,8,10,12\}\bigr)
\]
视为全部偶窗上的延拓候选。则纯计数必要条件已经强制
\[
2^m\ge F_{m+2}F_{m/2}.
\]
该条件在
\[
m=22
\]
时仍成立，但在
\[
m=24
\]
时首次失效；更精确地，
\[
F_{24}F_{11}=46368\cdot 89=4126752<2^{22}=4194304,
\]
而
\[
F_{26}F_{12}=121393\cdot 144=17480592>2^{24}=16777216.
\]
因此，若试图把审计偶窗上的最小退化 Fibonacci 律无条件外推到全部偶窗，则其首个纯鸽巢截断最迟发生在 \(m=24\)。
\end{theorem}
\begin{proof}
若对某个偶数 \(m\) 仍有
\[
d_{\min}(m)=F_{m/2},
\]
则每条纤维都至少含有 \(F_{m/2}\) 个微态。由稳定类型计数律 \(|X_m|=F_{m+2}\) 与恒等式
\[
2^m=\sum_{x\in X_m}d_m(x),
\]
立即得到必要条件
\[
2^m\ge |X_m|\,d_{\min}(m)=F_{m+2}F_{m/2}.
\]
对 \(m=22\) 与 \(m=24\) 直接代入，便得到
\[
F_{24}F_{11}=46368\cdot 89=4126752<2^{22}=4194304,
\]
\[
F_{26}F_{12}=121393\cdot 144=17480592>2^{24}=16777216.
\]
因此这一纯计数必要条件在 \(m=22\) 仍未破裂，而在 \(m=24\) 已经失效；于是任何无条件外推都不可能延续到 \(m=24\) 之后。证毕。
\end{proof}

\noindent
由此，最大非可缩纤维的模 \(6\) 相图、半阶精确暗块、window-\(6\) 根字典对 boundary 奇偶的三维盲核，以及黄金分支在 Fibonacci 分母上的双相输运常数，便不再构成彼此平行的孤立事实：前两者说明同一同余相上可以同时出现大规模谱消失与无折损拓扑主纤维，中者表明可见的根字典天然遗漏一条规范的中心奇偶方向，而纯鸽巢截断又说明最小退化 Fibonacci 律不能脱离审计窗口作无条件延展。因此，该折叠系统的可见图样同时受同余类拓扑、精确暗块与侧向采样异常三类彼此不可约机制共同支配。

\endinput

\paragraph{二原子经验谱、临界容量反演、Bernoulli 残差比与奇素数缺陷尾谱}
在 bin-fold 的两态一致渐近、critical-scale 容量极限曲线、fold-gauge 常数的 Bernoulli 递归剥离以及输出位势的奇素数单位根扭曲四阶展开已经建立之后，还可继续抽取出下列七条后果。它们分别刻画稳定层均匀经验谱的严格两点弱极限、连续容量极限曲线的分布二阶导数、目标成功率下的最优辅助比特预算、由 \((\log C_m)\) 序列恢复 Bernoulli 塔时的 Richardson 残差比、奇素数模 RH 缺陷的可和性与其 \(\ell^\alpha\) 临界指数，以及最小非平凡 prime-twist 谱半径对 \(\lambda\) 的严格二次接近律与最终单调逼近。

\begin{theorem}[bin-fold 纤维谱经自然缩放后的稳定层两点弱极限]\label{thm:xi-time-part70-binfold-uniform-empirical-two-point-law}
\leanverified{paper\_xi\_time\_part70\_binfold\_uniform\_empirical\_two\_point\_law}
定义尺度化多重度
\[
Z_m^{\mathrm{bin}}(w):=\left(\frac{\varphi}{2}\right)^m d_m^{\mathrm{bin}}(w)
\qquad (w\in X_m),
\]
以及稳定层上的均匀经验测度
\[
\nu_m:=\frac{1}{|X_m|}\sum_{w\in X_m}\delta_{Z_m^{\mathrm{bin}}(w)}.
\]
再记
\[
A_0^{(m)}:=\{w\in X_m:\ w_m=0\},
\qquad
A_1^{(m)}:=\{w\in X_m:\ w_m=1\}.
\]
则有弱收敛
\[
\nu_m\xrightarrow{\mathrm w}\nu_\infty
:=
\frac{1}{\varphi}\,\delta_1+\frac{1}{\varphi^2}\,\delta_{\varphi^{-1}}.
\]
等价地，对任意有界连续函数 \(f\)，都有
\[
\lim_{m\to\infty}\frac{1}{|X_m|}\sum_{w\in X_m}f\!\bigl(Z_m^{\mathrm{bin}}(w)\bigr)
=
\frac{1}{\varphi}f(1)+\frac{1}{\varphi^2}f(\varphi^{-1}).
\]
\end{theorem}
\begin{proof}
由定理 \ref{thm:fold-bin-two-state-asymptotic}，存在 \(\varepsilon_m\to 0\) 使得对全部 \(w\in X_m\) 一致有
\[
\bigl|Z_m^{\mathrm{bin}}(w)-\varphi^{-w_m}\bigr|\le \varepsilon_m.
\]
因此对充分大的 \(m\)，所有 \(Z_m^{\mathrm{bin}}(w)\) 都落在某个固定紧区间内。任取有界连续函数 \(f\)，其在该紧区间上一致连续。于是若定义
\[
\eta_{m,0}:=\sup_{w\in A_0^{(m)}}\bigl|f(Z_m^{\mathrm{bin}}(w))-f(1)\bigr|,
\qquad
\eta_{m,1}:=\sup_{w\in A_1^{(m)}}\bigl|f(Z_m^{\mathrm{bin}}(w))-f(\varphi^{-1})\bigr|,
\]
则 \(\eta_{m,0},\eta_{m,1}\to 0\)。从而
\[
\frac{1}{|X_m|}\sum_{w\in X_m}f\!\bigl(Z_m^{\mathrm{bin}}(w)\bigr)
=
\frac{|A_0^{(m)}|}{|X_m|}\bigl(f(1)+o(1)\bigr)
+
\frac{|A_1^{(m)}|}{|X_m|}\bigl(f(\varphi^{-1})+o(1)\bigr).
\]
再由稳定词计数
\[
|A_0^{(m)}|=F_{m+1},
\qquad
|A_1^{(m)}|=F_m,
\qquad
|X_m|=F_{m+2},
\]
以及 Fibonacci 比值极限
\[
\frac{F_{m+1}}{F_{m+2}}\to \frac{1}{\varphi},
\qquad
\frac{F_m}{F_{m+2}}\to \frac{1}{\varphi^2},
\]
便得到所示极限。证毕。
\end{proof}

\begin{theorem}[连续容量极限曲线的分布二阶导数恰为两原子测度]\label{thm:xi-time-part70-continuous-capacity-two-atom-second-derivative}
定义尺度化连续容量
\[
\widehat C_m(s):=
2^{-m}\,\mathcal C_m^{\mathrm{bin,cont}}\!\left(
s\left(\frac{2}{\varphi}\right)^m
\right)
\qquad (s\ge 0),
\]
并记其极限函数
\[
L(s):=\frac{\varphi}{\sqrt5}\min\{1,s\}+\frac{1}{\sqrt5}\min\{\varphi^{-1},s\}.
\]
则在分布意义下有精确恒等式
\[
-D^2L
=
\frac{1}{\sqrt5}\,\delta_{\varphi^{-1}}
+
\frac{\varphi}{\sqrt5}\,\delta_1.
\]
换言之，极限容量曲线只有两个不可消去的折点，其位置分别位于 \(s=\varphi^{-1}\) 与 \(s=1\)，对应原子质量分别为 \(1/\sqrt5\) 与 \(\varphi/\sqrt5\)。
\end{theorem}
\begin{proof}
这正是定理 \ref{thm:xi-foldbin-underresolution-two-atomic-curvature-measure} 中曲率测度公式的直接重写。若将
\[
L(s)=
\begin{cases}
\dfrac{\varphi^2}{\sqrt5}\,s, & 0\le s\le \varphi^{-1},\\[6pt]
\dfrac{\varphi}{\sqrt5}\,s+\dfrac{1}{\varphi\sqrt5}, & \varphi^{-1}\le s\le 1,\\[8pt]
1, & s\ge 1
\end{cases}
\]
逐段求导，则 \(L'\) 在 \(s=\varphi^{-1}\) 处的跃降为
\[
\frac{\varphi^2}{\sqrt5}-\frac{\varphi}{\sqrt5}=\frac{1}{\sqrt5},
\]
在 \(s=1\) 处的跃降为
\[
\frac{\varphi}{\sqrt5}-0=\frac{\varphi}{\sqrt5}.
\]
因此分布导数满足
\[
-D^2L
=
\frac{1}{\sqrt5}\,\delta_{\varphi^{-1}}
+
\frac{\varphi}{\sqrt5}\,\delta_1.
\]
证毕。
\end{proof}

\begin{theorem}[目标成功率下最优辅助比特预算的严格分段阈值]\label{thm:xi-time-part70-target-success-optimal-budget-threshold}
沿用定理 \ref{thm:xi-foldbin-dyadic-capacity-critical-window-limit} 的记号。对任意目标成功率 \(\sigma\in(0,1)\)，定义最小预算
\[
B_m^\star(\sigma):=
\min\bigl\{B\in\mathbb N:\ \mathrm{Succ}_m^{\mathrm{bin}}(B)\ge \sigma\bigr\}.
\]
再定义连续阈值常数
\[
s_\sigma:=
\begin{cases}
\dfrac{\sqrt5}{\varphi^2}\,\sigma,
& 0<\sigma\le \dfrac{\varphi}{\sqrt5},\\[8pt]
\dfrac{\sqrt5\,\sigma-\varphi^{-1}}{\varphi},
& \dfrac{\varphi}{\sqrt5}\le \sigma<1.
\end{cases}
\]
则有
\[
B_m^\star(\sigma)
=
m\log_2\!\left(\frac{2}{\varphi}\right)+\log_2 s_\sigma+O(1).
\]
更一般地，若整数序列 \((B_m)_{m\ge 1}\) 满足
\[
2^{B_m}\left(\frac{\varphi}{2}\right)^m\longrightarrow s\in(0,\infty),
\]
则
\[
\mathrm{Succ}_m^{\mathrm{bin}}(B_m)\longrightarrow L(s),
\]
其中 \(L\) 即定理 \ref{thm:xi-time-part70-continuous-capacity-two-atom-second-derivative} 的极限曲线；因此 \(s_\sigma\) 正是方程 \(L(s)=\sigma\) 的唯一解。
\end{theorem}
\begin{proof}
第二个结论正是定理 \ref{thm:xi-foldbin-dyadic-capacity-critical-window-limit} 与推论 \ref{cor:xi-foldbin-critical-scale-success-curve-inversion} 的内容，其中极限成功曲线即
\[
L(s)=
\begin{cases}
\dfrac{\varphi^2}{\sqrt5}\,s, & 0\le s\le \varphi^{-1},\\[6pt]
\dfrac{\varphi}{\sqrt5}\,s+\dfrac{1}{\varphi\sqrt5}, & \varphi^{-1}\le s\le 1,\\[8pt]
1, & s\ge 1.
\end{cases}
\]
因此 \(s_\sigma=L^{-1}(\sigma)\) 的显式公式恰为上式两段线性方程的解。

现证明 \(B_m^\star(\sigma)\) 的渐近公式。由于 \(\sigma\in(0,1)\)，故 \(s_\sigma\in(0,1)\)。取
\[
0<\varepsilon<\min\{s_\sigma,1-s_\sigma\}.
\]
由 \(L\) 在 \([0,1]\) 上严格递增可知
\[
L(s_\sigma-\varepsilon)<\sigma<L(s_\sigma+\varepsilon).
\]
再由 critical-scale 极限，分别取
\[
\widetilde B_m^-:=\left\lfloor \log_2\!\Bigl((s_\sigma-\varepsilon)\Bigl(\frac{2}{\varphi}\Bigr)^m\Bigr)\right\rfloor,
\qquad
\widetilde B_m^+:=\left\lceil \log_2\!\Bigl((s_\sigma+\varepsilon)\Bigl(\frac{2}{\varphi}\Bigr)^m\Bigr)\right\rceil,
\]
则当 \(m\) 充分大时有
\[
\mathrm{Succ}_m^{\mathrm{bin}}(\widetilde B_m^-)<\sigma,
\qquad
\mathrm{Succ}_m^{\mathrm{bin}}(\widetilde B_m^+)\ge \sigma.
\]
由于 \(\mathrm{Succ}_m^{\mathrm{bin}}(B)\) 关于 \(B\) 单调不减，而 \(B_m^\star(\sigma)\) 按定义是达到成功率阈值的最小整数预算，遂得
\[
\widetilde B_m^-+1\le B_m^\star(\sigma)\le \widetilde B_m^+.
\]
把上式展开后便得到
\[
\log_2\!\Bigl((s_\sigma-\varepsilon)\Bigl(\frac{2}{\varphi}\Bigr)^m\Bigr)
\le
B_m^\star(\sigma)
\le
\log_2\!\Bigl((s_\sigma+\varepsilon)\Bigl(\frac{2}{\varphi}\Bigr)^m\Bigr)+1.
\]
由于 \(\varepsilon\) 固定，端点与
\[
m\log_2\!\left(\frac{2}{\varphi}\right)+\log_2 s_\sigma
\]
之差都是有界量，故结论成立。证毕。
\end{proof}

\begin{theorem}[由 \texorpdfstring{$(\log C_m)$}{(log Cm)} 序列恢复 Bernoulli 塔时的 Richardson 残差比极限]\label{thm:xi-time-part70-logcm-richardson-residual-ratio}
记
\[
q:=\frac{\varphi}{2}\in(0,1),
\qquad
E_m:=\log C_m-\log(2\pi),
\]
\[
a_r:=
\frac{B_{2r}}{r(2r-1)}
\bigl(\varphi^{-1}+\varphi^{2r-3}\bigr)
\qquad (r\ge 1).
\]
对每个固定整数 \(R\ge 0\)，定义 \(R\) 阶 Richardson 残差
\[
\mathcal R_m^{(R)}
:=
E_m-\sum_{r=1}^{R}a_r q^{(2r-1)m}.
\]
则有严格比值极限
\[
\frac{\mathcal R_{m+1}^{(R)}}{\mathcal R_m^{(R)}}
\longrightarrow
q^{2R+1}
=
\left(\frac{\varphi}{2}\right)^{2R+1}.
\]
换言之，一旦减去前 \(R\) 层 Bernoulli 模态，剩余误差的相邻比值必然锁定到奇数次黄金幂 \((\varphi/2)^{2R+1}\)。
\end{theorem}
\begin{proof}
对任意固定 \(R\ge 0\)，将定理 \ref{thm:fold-bin-gauge-constant-stirling-bernoulli-hierarchy} 应用于 \(R+1\) 阶，得到
\[
E_m
=
\sum_{r=1}^{R+1}a_r q^{(2r-1)m}
+O\!\bigl(q^{(2R+3)m}\bigr).
\]
因此
\[
\mathcal R_m^{(R)}
=
a_{R+1}q^{(2R+1)m}
+O\!\bigl(q^{(2R+3)m}\bigr).
\]
又因为偶 Bernoulli 数满足 \(B_{2r}\neq 0\)（\(r\ge 1\)），而 \(\varphi^{-1}+\varphi^{2r-3}>0\)，故
\[
a_{R+1}\neq 0.
\]
于是
\[
\mathcal R_m^{(R)}
=
a_{R+1}q^{(2R+1)m}\bigl(1+O(q^{2m})\bigr),
\]
从而
\[
\frac{\mathcal R_{m+1}^{(R)}}{\mathcal R_m^{(R)}}
=
\frac{a_{R+1}q^{(2R+1)(m+1)}\bigl(1+O(q^{2m})\bigr)}
{a_{R+1}q^{(2R+1)m}\bigl(1+O(q^{2m})\bigr)}
\longrightarrow
q^{2R+1}.
\]
证毕。
\end{proof}

\begin{theorem}[奇素数模 RH 缺陷总质量的可和性]\label{thm:xi-time-part70-odd-prime-rh-defect-summable-mass}
记
\[
\lambda:=\varphi^2,
\qquad
\rho_{p,1}:=\rho\!\bigl(M(e^{2\pi i/p})\bigr)
\qquad
(p\ge 3\ \text{为奇素数}),
\]
并定义扭曲指数与其缺陷
\[
\theta_{p,1}:=\frac{\log \rho_{p,1}}{\log\lambda},
\qquad
\delta_p:=1-\theta_{p,1}.
\]
则有显式展开
\[
\delta_p
=
\frac{4\pi^2(6\sqrt5-5)}{250\log(\varphi^2)}\,\frac{1}{p^2}
+O\!\left(\frac{1}{p^4}\right).
\]
从而
\[
\sum_{\substack{p\ge 3\\ p\ \text{为奇素数}}}\delta_p<\infty,
\qquad
\sum_{\substack{p>P\\ p\ \text{为奇素数}}}\delta_p=O(P^{-1}).
\]
因此，尽管奇素数方向上的平方根门槛失效在频率上呈余有限出现，其缺陷幅度却形成可和尾谱。
\end{theorem}
\begin{proof}
由定理 \ref{thm:xi-output-potential-odd-prime-rh-obstruction-index-to-one}，
\[
\theta_{p,1}
=
1-\frac{6\sqrt5-5}{250\log\lambda}\Bigl(\frac{2\pi}{p}\Bigr)^2
+\frac{7+24\sqrt5}{75000\log\lambda}\Bigl(\frac{2\pi}{p}\Bigr)^4
+O(p^{-6}).
\]
将 \(\lambda=\varphi^2\) 代入并移项，便得到
\[
\delta_p
=
\frac{4\pi^2(6\sqrt5-5)}{250\log(\varphi^2)}\,\frac{1}{p^2}
+O(p^{-4}).
\]
特别地，存在常数 \(c_1,c_2>0\) 与奇素数阈值 \(p_0\)，使得对所有奇素数 \(p\ge p_0\) 都有
\[
\frac{c_1}{p^2}\le \delta_p\le \frac{c_2}{p^2}.
\]
于是
\[
\sum_{\substack{p\ge p_0\\ p\ \text{为奇素数}}}\delta_p
\ll
\sum_{n\ge p_0}\frac{1}{n^2}
<\infty,
\]
从而总和收敛。对尾和，同样有
\[
\sum_{\substack{p>P\\ p\ \text{为奇素数}}}\delta_p
\ll
\sum_{n>P}\frac{1}{n^2}
=
O(P^{-1}).
\]
证毕。
\end{proof}

\begin{corollary}[奇素数 RH 缺陷尾谱的 \texorpdfstring{$\ell^\alpha$}{ell alpha} 临界指数恰为 \texorpdfstring{$1/2$}{1/2}]\label{cor:xi-time-part70-odd-prime-rh-defect-ell-alpha-threshold}
沿用定理 \ref{thm:xi-time-part70-odd-prime-rh-defect-summable-mass} 的记号。则对任意 \(\alpha>0\)，有
\[
(\delta_p)_{p\ \text{为奇素数}}\in \ell^\alpha
\qquad\Longleftrightarrow\qquad
\alpha>\frac12.
\]
\end{corollary}
\begin{proof}
由定理 \ref{thm:xi-time-part70-odd-prime-rh-defect-summable-mass}，存在常数 \(c_1,c_2>0\) 与 \(p_0\) 使得对全部奇素数 \(p\ge p_0\) 都有
\[
c_1p^{-2}\le \delta_p\le c_2p^{-2}.
\]
于是
\[
\delta_p^\alpha\asymp p^{-2\alpha}
\qquad (p\to\infty,\ p\ \text{为奇素数}).
\]
而素数级数满足
\[
\sum_{p}p^{-s}<\infty
\qquad\Longleftrightarrow\qquad
s>1.
\]
因此
\[
\sum_{\substack{p\ge 3\\ p\ \text{为奇素数}}}\delta_p^\alpha<\infty
\qquad\Longleftrightarrow\qquad
\sum_{\substack{p\ge 3\\ p\ \text{为奇素数}}}p^{-2\alpha}<\infty
\qquad\Longleftrightarrow\qquad
2\alpha>1,
\]
即所求结论。证毕。
\end{proof}

\begin{theorem}[奇素数最小非平凡单位根扭曲的谱半径二次接近律与最终严格单调逼近]\label{thm:xi-time-part70-prime-twist-spectral-radius-quadratic-approach}
沿用推论 \ref{cor:real-input-40-collision-twist-fourth} 的记号，并记
\[
\lambda:=\rho(1)=\varphi^2,
\qquad
\rho_p:=\rho\!\left(e^{2\pi i/p}\right)
\qquad
(p\ge 3\ \text{为奇素数}).
\]
则成立：
\begin{enumerate}
  \item 有精确二次接近律
  \[
  \boxed{\;
  p^2\left(1-\frac{\rho_p}{\lambda}\right)
  \longrightarrow
  \frac{6\sqrt5-5}{250}\,(2\pi)^2.\;}
  \]
  \item 存在奇素数阈值 \(p_0\)，使得对任意奇素数 \(p<q\) 且 \(p,q\ge p_0\)，都有
  \[
  \boxed{\;
  \rho_p<\rho_q<\lambda.\;}
  \]
\end{enumerate}
换言之，最小非平凡 prime twist 的谱半径不仅以严格 \(p^{-2}\) 主项从下方逼近 \(\lambda\)，而且在充分大素数区间上按模数严格单调恢复。
\end{theorem}
\begin{proof}
由推论 \ref{cor:real-input-40-collision-twist-fourth}，当 \(t\to 0\) 时有
\[
\log\frac{\rho(t)}{\lambda}
=
-A t^2+B t^4+O(t^6),
\]
其中
\[
A:=\frac{6\sqrt5-5}{250}>0,
\qquad
B:=\frac{7+24\sqrt5}{75000}.
\]
将其指数化，得到
\[
\frac{\rho(t)}{\lambda}
=
\exp\!\bigl(-A t^2+B t^4+O(t^6)\bigr)
=
1-A t^2+O(t^4).
\]
对 \(t=t_p:=2\pi/p\) 代入，即得
\[
1-\frac{\rho_p}{\lambda}
=
A\left(\frac{2\pi}{p}\right)^2+O(p^{-4}),
\]
从而
\[
p^2\left(1-\frac{\rho_p}{\lambda}\right)
\longrightarrow
A(2\pi)^2
=
\frac{6\sqrt5-5}{250}\,(2\pi)^2,
\]
这就证明了第一条。

再令
\[
R(t):=\frac{\rho(t)}{\lambda}.
\]
由上式可得
\[
\log R(t)=-A t^2+B t^4+O(t^6).
\]
对其求导，得到
\[
\frac{R'(t)}{R(t)}
=
-2At+4Bt^3+O(t^5)
=
-2At\bigl(1+O(t^2)\bigr).
\]
因此存在 \(t_0>0\)，使得当 \(0<t\le t_0\) 时，
\[
\frac{R'(t)}{R(t)}<0.
\]
又因 \(R(t)>0\)，可知 \(R'(t)<0\)；故 \(R\) 在 \((0,t_0]\) 上严格递减。另一方面，由
\[
R(t)=1-A t^2+O(t^4)
\]
可知在同一小邻域内 \(R(t)<1\)。

现取奇素数阈值 \(p_0\) 使得
\[
\frac{2\pi}{p_0}\le t_0.
\]
若 \(p<q\) 且 \(p,q\ge p_0\)，则
\[
t_p=\frac{2\pi}{p}>\frac{2\pi}{q}=t_q,
\]
而 \(R\) 在 \((0,t_0]\) 上严格递减，遂得
\[
R(t_p)<R(t_q)<1.
\]
两边同乘以 \(\lambda\)，即
\[
\rho_p<\rho_q<\lambda.
\]
第二条得证。证毕。
\end{proof}

\noindent
上述七条结果表明，二原子塌缩并不只停留在 bin-fold 的输出实验或连续容量曲线上，而会同时外显为稳定层均匀经验谱的严格弱极限、目标成功率反演中的唯一分段阈值、\((\log C_m)\) 的奇次黄金幂残差比，以及奇素数单位根扭曲缺陷的可和尾谱与谱半径恢复律。因而，\(\{\varphi^{-1},1\}\) 这两个黄金比例尺度既支配欠分辨容量几何，也支配 Bernoulli 递归剥离后的剩余模态；而在奇素数单位根扭曲方向上，坏模数的频率与缺陷的振幅不仅被严格分离为“余有限出现”与“总质量可和”两种不同层级，而且最小非平凡扭曲的谱半径还会按严格 \(p^{-2}\) 主项最终单调地从下方逼近 \(\lambda\)。于是，稳定层二点律、资源阈值反演、Richardson 诊断与奇素数 RH 缺陷最终被压缩进同一条黄金尺度主线。

\input{subsubsec__xi-time-protocol-conclusions-part70a-uniform-twostate-gauge-firstorder}

\input{subsubsec__xi-time-protocol-conclusions-part70aa-real-moment-inversemoment-bernoulli-bridge}

\input{subsubsec__xi-time-protocol-conclusions-part70ab-terminalbit-threshold-entropy-mellin-bridge}
\input{subsubsec__xi-time-protocol-conclusions-part70ac-hologram-prefix-entropy-minbit}
\input{subsubsec__xi-time-protocol-conclusions-part70ad-binfold-information-geometry-renyi-tsallis-gauge}

\input{subsubsec__xi-time-protocol-conclusions-part70b-double-kink-escort-tail-darkband-interface}

\input{subsubsec__xi-time-protocol-conclusions-part70c-spectrum-z2-orbit-cancellation}

\input{subsubsec__xi-time-protocol-conclusions-part70c-zeroorbit-helly-nerve-density-kl}

\input{subsubsec__xi-time-protocol-conclusions-part70d-terminalsheet-rigidity-critical-oscillation-riesz-excess}
\input{subsubsec__xi-time-protocol-conclusions-part70da-window6-orbittype-c12-geometric-obstruction}
\input{subsubsec__xi-time-protocol-conclusions-part70e-binfold-entire-amplitude-stokes-rigidity}

\input{subsubsec__xi-time-protocol-conclusions-part71-zero-coset-clique-resonance-coefficients}

\endinput


\endinput


\endinput

\paragraph{二步 primitive 全系数、偶模坏模律与误差指数的反向成本}
在七条分支二回路的 primitive 塌缩、模 \(2\) 的最小坏模证人、正实退化点列表以及纯碰撞误差指数的相变已经建立之后，还可进一步抽取出下列五条完全闭式的后果。它们分别把最短 primitive 层的全系数、其在 ghost 层的偶次迹塔、偶数模上的坏模普遍性、\(\theta=0\) 邻域无退化窗的极大性，以及 \(\beta(u)\to 1\) 的前向渐近所对应的反向参数成本压缩为显式公式。

\begin{theorem}[真实输入输出位势的二步 primitive 系数恰为 \texorpdfstring{$8u$}{8u}]\label{thm:xi-time-part68a-output-two-step-primitive-eight-u}
\leanverified{paper\_xi\_time\_part68a\_output\_two\_step\_primitive\_eight\_u}
沿用定理 \ref{thm:xi-time-part62d-atomic-two-step-witt-rankone-collapse} 的记号，记
\[
P(z,u)=\sum_{n\ge 1}p_n(u)z^n.
\]
则二步 primitive 系数满足精确恒等式
\[
p_2(u)=8u.
\]
\end{theorem}
\begin{proof}
由定理 \ref{thm:xi-time-part62d-atomic-two-step-witt-rankone-collapse}，
\[
p_2(u)=u+p_2^{\mathrm{core}}(u).
\]
同一定理并已指出：branch 层中满足
\[
(|\tau|,E)=(2,1)
\]
的两步回路恰有 \(7\) 条。按 primitive 生成函数的定义，每一条长度 \(2\)、能量 \(1\) 的 primitive 回路都恰贡献一个 monomial \(u z^2\)；因此
\[
p_2^{\mathrm{core}}(u)=7u.
\]
代回即得
\[
p_2(u)=u+7u=8u.
\]
证毕。
\end{proof}

\begin{corollary}[长度 \texorpdfstring{$2$}{2} primitive 扇区在 ghost 层给出完全显式的偶次迹塔]\label{cor:xi-time-part68a-output-two-step-primitive-even-ghost-tower}
沿用定理 \ref{thm:xi-time-part68a-output-two-step-primitive-eight-u} 的记号。定义由长度 \(2\) primitive 扇区单独诱导的 ghost 级数为
\[
\sum_{n\ge 1}\frac{a_n^{\langle 2\rangle}(u)}{n}z^n
:=
\sum_{m\ge 1}\frac{p_2(u^m)}{m}z^{2m}.
\]
则对一切 \(m\ge 1\) 都有
\[
a_{2m}^{\langle 2\rangle}(u)=16u^m,
\qquad
a_{2m+1}^{\langle 2\rangle}(u)=0.
\]
\end{corollary}
\begin{proof}
由定理 \ref{thm:xi-time-part68a-output-two-step-primitive-eight-u}，
\[
p_2(u^m)=8u^m
\qquad (m\ge 1).
\]
故
\[
\sum_{n\ge 1}\frac{a_n^{\langle 2\rangle}(u)}{n}z^n
=
\sum_{m\ge 1}\frac{8u^m}{m}z^{2m}.
\]
将右端与
\[
\sum_{m\ge 1}\frac{a_{2m}^{\langle 2\rangle}(u)}{2m}z^{2m}
\]
逐项比较，立得
\[
\frac{a_{2m}^{\langle 2\rangle}(u)}{2m}=\frac{8u^m}{m},
\]
从而
\[
a_{2m}^{\langle 2\rangle}(u)=16u^m.
\]
而右端仅含偶次幂，故全部奇次系数必为 \(0\)。证毕。
\end{proof}

\begin{theorem}[每一个偶数模都是坏模数]\label{thm:xi-time-part68a-collision-all-even-moduli-bad}
沿用定理 \ref{thm:xi-real-input-40-collision-minimal-modulus-two-witness} 的记号。对每个模数 \(m\ge 2\) 与扭曲通道 \(1\le j\le m-1\)，记
\[
\omega_m:=e^{2\pi i/m},
\qquad
\rho_{m,j}:=\rho\!\bigl(W(\omega_m^j)\bigr),
\qquad
\lambda:=\rho(W(1))=\varphi^2.
\]
则对每一个偶数 \(m\) 都存在某个 \(j\) 使
\[
\rho_{m,j}>\sqrt{\lambda}=\varphi.
\]
更具体地，取 \(j=m/2\) 即有
\[
\rho_{m,m/2}=\rho(W(-1))>\varphi.
\]
因此全部偶数模都是平方根门槛失效的坏模数。
\end{theorem}
\begin{proof}
若 \(m\) 为偶数，则取
\[
j=\frac{m}{2}.
\]
此时
\[
\omega_m^{\,m/2}=e^{\pi i}=-1,
\]
故
\[
\rho_{m,m/2}=\rho\!\bigl(W(-1)\bigr).
\]
而定理 \ref{thm:xi-real-input-40-collision-minimal-modulus-two-witness} 已证明
\[
\rho\!\bigl(W(-1)\bigr)>\sqrt{\lambda}=\varphi.
\]
于是该偶数模 \(m\) 至少在通道 \(j=m/2\) 上突破平方根门槛。证毕。
\end{proof}

\begin{theorem}[\texorpdfstring{$\theta=0$}{theta=0} 邻域正实参数的最大对称无退化窗]\label{thm:xi-time-part68a-output-positive-real-maximal-nondegenerate-window}
沿用注 \ref{rem:real-input-40-output-potential-degeneracy-phase} 的记号。记距离 \(u=1\) 最近的两个正实退化点为
\[
u_-=0.9328372624859649\ldots,
\qquad
u_+=1.0547964213679192\ldots,
\]
并定义
\[
\theta_0:=\min\{-\log u_-,\,\log u_+\}
=0.0533477827\ldots.
\]
则：
\begin{enumerate}
  \item 当 \(|\theta|<\theta_0\) 时，沿正实参数 \(u=e^\theta\) 不会跨越任何正实退化点；
  \item 若 \(\theta_1>\theta_0\)，则对称区间 \((-\theta_1,\theta_1)\) 必跨越至少一个正实退化点。
\end{enumerate}
因而 \((-\theta_0,\theta_0)\) 正是以 \(\theta=0\) 为中心的最大对称无退化实窗。
\end{theorem}
\begin{proof}
注 \ref{rem:real-input-40-output-potential-degeneracy-phase} 已给出全部正实退化点
\[
u\in\{0.008460\ldots,\ 0.745534\ldots,\ 0.932837\ldots,\ 1.054796\ldots,\ 2.090392\ldots,\ 15.522820\ldots\},
\]
并指出离 \(u=1\) 最近的两点正是 \(u_-\) 与 \(u_+\)。

若 \(|\theta|<\theta_0\)，则
\[
-\log u_-<\theta<\log u_+,
\]
从而
\[
u_-<e^\theta<u_+.
\]
因此 \(u=e^\theta\) 落在开区间 \((u_-,u_+)\) 内，不会跨越任何正实退化点。这证明了第 \(1\) 条。

反之，若 \(\theta_1>\theta_0\)，则由 \(\theta_0\) 的定义可知，必有
\[
\theta_1>-\log u_-
\quad\text{或}\quad
\theta_1>\log u_+.
\]
在前一种情形，区间 \((-\theta_1,\theta_1)\) 包含点 \(\log u_-<0\)；在后一种情形，它包含点 \(\log u_+>0\)。两者都意味着该对称区间至少跨越一个正实退化点。故第 \(2\) 条成立，极大性随之得到。证毕。
\end{proof}

\begin{theorem}[纯碰撞误差指数的 Lambert--\texorpdfstring{$\mathscr W$}{W} 反演律]\label{thm:xi-time-part68a-collision-beta-lambertw-inversion}
沿用推论 \ref{cor:real-input-40-collision-error-phase} 的记号，并令
\[
\varepsilon:=1-\beta(u)\downarrow 0
\qquad (u\to+\infty).
\]
记 \(\mathscr W\) 为 Lambert--\(\mathscr W\) 函数主支，即
\[
\mathscr W(x)e^{\mathscr W(x)}=x
\qquad (x>0).
\]
则有反演渐近
\[
\sqrt{u}
=
\frac{1}{\varepsilon\,\mathscr W(1/\varepsilon)}(1+o(1)),
\qquad
u
=
\frac{1}{\varepsilon^2\,\mathscr W(1/\varepsilon)^2}(1+o(1)).
\]
特别地，由
\[
\mathscr W(x)=\log x-\log\log x+o(1)
\qquad (x\to+\infty)
\]
可得粗尺度资源律
\[
u\asymp \frac{1}{\varepsilon^2\log^2(1/\varepsilon)}.
\]
\end{theorem}
\begin{proof}
由推论 \ref{cor:real-input-40-collision-error-phase}，
\[
\beta(u)=1-\frac{2}{\sqrt u\,\log u}+O\!\left(\frac{1}{u\log u}\right)
\qquad (u\to+\infty).
\]
令
\[
x:=\sqrt u.
\]
则 \(\log u=2\log x\)，故
\[
\varepsilon
=
\frac{1}{x\log x}+O\!\left(\frac{1}{x^2\log x}\right)
=
\frac{1}{x\log x}(1+o(1)).
\]
于是
\[
x\log x=\frac{1}{\varepsilon}(1+o(1)).
\]
再令
\[
y:=\log x,
\]
则 \(x=e^y\)，从而
\[
y e^y=\frac{1}{\varepsilon}(1+o(1)).
\]
按 Lambert--\(\mathscr W\) 函数定义，
\[
y=\mathscr W\!\left(\frac{1}{\varepsilon}(1+o(1))\right).
\]
又由于
\[
\mathscr W'(t)=\frac{\mathscr W(t)}{t(1+\mathscr W(t))}=O(t^{-1})
\qquad (t\to+\infty),
\]
故
\[
\mathscr W\!\left(\frac{1}{\varepsilon}(1+o(1))\right)
=
\mathscr W\!\left(\frac{1}{\varepsilon}\right)+o(1).
\]
于是
\[
x=e^y
=
\frac{(1/\varepsilon)(1+o(1))}{\mathscr W(1/\varepsilon)+o(1)}
=
\frac{1}{\varepsilon\,\mathscr W(1/\varepsilon)}(1+o(1)).
\]
这就得到
\[
\sqrt u=\frac{1}{\varepsilon\,\mathscr W(1/\varepsilon)}(1+o(1)).
\]
两边平方即得
\[
u=\frac{1}{\varepsilon^2\,\mathscr W(1/\varepsilon)^2}(1+o(1)).
\]
最后由
\[
\mathscr W(1/\varepsilon)\sim \log(1/\varepsilon)
\qquad (\varepsilon\downarrow 0)
\]
可知
\[
u\asymp \frac{1}{\varepsilon^2\log^2(1/\varepsilon)}.
\]
证毕。
\end{proof}

\noindent
由此，真实输入 \(40\) 状态核的最短 primitive 层、偶模单位根扭曲、正实退化窗与纯碰撞误差指数之间形成了一条更细的封闭链：二步 primitive 坐标已在 \(8u\) 处完全闭式化，并在 ghost 层外显为一整条偶次迹塔；奇偶扭曲迫使全部偶数模统一越过平方根门槛，而正实参数的中心对称无退化窗则恰由 \(\theta_0\) 取尽；最后，\(\beta(u)\to 1\) 的慢模塌缩不只给出前向渐近，而且可反解为带 Lambert--\(\mathscr W\) 校正的显式参数成本律。

\endinput

\paragraph{围绕 \texorpdfstring{$u=1$}{u=1} 的正实谱室与代数 \texorpdfstring{$\pm1$}{±1} 锚点}
在正实退化点列表、中心对称无退化窗以及整数级共振校准点都已确定之后，还可把真实输入 \(40\) 状态核的正实谱几何进一步压缩为下列两条定理。前一条把围绕 \(u=1\) 的最大正实谱室与其对数半径完全闭式化；后一条则把 \(\lambda=\pm1\) 的两枚正实代数锚点提升为非退化谱证书。

\begin{theorem}[围绕 \texorpdfstring{$u=1$}{u=1} 的正实最大无额外碰撞谱室]\label{thm:xi-time-part68aa-output-positive-real-maximal-collisionfree-chamber}
沿用命题 \ref{prop:real-input-40-output-spectral-collisions} 与注 \ref{rem:real-input-40-output-potential-degeneracy-phase} 的记号，记
\[
\Delta(z,u)=(1-uz^2)Q(z,u).
\]
则在正实参数轴 \(u>0\) 上，除结构性点 \(u=1\) 外，\(Q(\cdot,u)\) 出现重根当且仅当
\[
u\in
\left\{
0.0084602654608787966132,\ 
0.74553449580298956225,\ 
0.93283726248596492776,\ 
1.0547964213679192275,\ 
2.0903923780490870311,\ 
15.522820704426852031
\right\}.
\]
因此，以 \(u=1\) 为中心且不再跨越其他正实碰撞点的最大开区间恰为
\[
\bigl(0.93283726248596492776,\ 1.0547964213679192275\bigr).
\]
若改用
\[
\theta:=\log u
\]
作为局部参数，则以 \(\theta=0\) 为中心的最大对称无额外碰撞区间为
\[
|\theta|<\theta_0,
\qquad
\theta_0
:=
\min\!\Bigl\{
\log 1.0547964213679192275,\ 
-\log 0.93283726248596492776
\Bigr\}
=
0.0533477827\ldots.
\]
\end{theorem}
\begin{proof}
注 \ref{rem:real-input-40-output-potential-degeneracy-phase} 已列出 \(u>0\) 上全部正实退化点，并指出离 \(u=1\) 最近的两点恰为
\[
u_-=0.93283726248596492776,\qquad
u_+=1.0547964213679192275.
\]
故在去掉这些碰撞点之后，包含 \(u=1\) 的最大连通开区间只能是
\[
(u_-,u_+)
=
\bigl(0.93283726248596492776,\ 1.0547964213679192275\bigr).
\]
这就给出了第一条结论。

再取 \(\theta=\log u\)。由于对数映射在 \((0,+\infty)\) 上严格单调，上一开区间对应于
\[
\bigl(\log u_-,\ \log u_+\bigr).
\]
若要求得到以 \(\theta=0\) 为中心的对称区间，则其最大半径只能取左右两端到 \(0\) 的较小距离，即
\[
\theta_0=\min\{-\log u_-,\,\log u_+\}.
\]
代入上述数值便得
\[
\theta_0=0.0533477827\ldots.
\]
证毕。
\end{proof}

\begin{theorem}[非退化 \texorpdfstring{$\pm1$}{±1} 谱锚的代数精确性]\label{thm:xi-time-part68aa-output-algebraic-plusminus-one-nondegenerate-anchors}
\leanverified{paper\_xi\_time\_part68aa\_output\_algebraic\_plusminus\_one\_nondegenerate\_anchors}
沿用
\[
\Delta(z,u)=\det(I-zM(u))=(1-uz^2)Q(z,u)
\]
的记号。则存在两枚完全代数的正实参数
\[
u_+=4,
\qquad
u_-=-3+\sqrt{11},
\]
使得矩阵 \(M(u)\) 分别具有特征值
\[
\lambda=1
\qquad\text{与}\qquad
\lambda=-1.
\]
更强地，这两个谱锚都落在正实碰撞集合之外，因而对应特征值均为非退化谱点。换言之，\(\lambda=\pm1\) 在这两处都给出完全代数且可直接审计的正实谱校准点。
\end{theorem}
\begin{proof}
由推论 \ref{cor:real-input-40-integer-resonance-calibration}，
\[
Q(1,4)=0,
\qquad
Q(-1,-3+\sqrt{11})=0.
\]
当 \(u\neq 1\) 时，
\[
\Delta(\pm1,u)=(1-u)Q(\pm1,u),
\]
故
\[
\Delta(1,4)=0,
\qquad
\Delta(-1,-3+\sqrt{11})=0.
\]
按
\[
\Delta(z,u)=\det(I-zM(u)),
\]
可知 \(z=\pm1\) 为零点等价于 \(M(u)\) 具有特征值 \(\lambda=1/z=\pm1\)。因此在 \(u_+=4\) 与 \(u_-=-3+\sqrt{11}\) 处分别得到 \(\lambda=1\) 与 \(\lambda=-1\)。

另一方面，注 \ref{rem:real-input-40-output-potential-degeneracy-phase} 已给出全部正实碰撞参数。显然
\[
4\notin
\left\{
0.008460\ldots,\ 0.745534\ldots,\ 0.932837\ldots,\ 1.054796\ldots,\ 2.090392\ldots,\ 15.522820\ldots
\right\},
\]
且
\[
-3+\sqrt{11}\approx 0.31662479
\]
同样不在该集合中，并且二者都不等于结构性点 \(u=1\)。因此对应零点都是简单根，亦即 \(\lambda=\pm1\) 在这两处都是非退化谱点。证毕。
\end{proof}

\noindent
由此，围绕平衡点 \(u=1\) 的正实谱室不再只由局部展开半径间接刻画，而是被一张完整的碰撞点表和一个最优对数半径精确锁定；与此同时，\(\lambda=\pm1\) 的两枚正实代数锚点又把整数级共振提升为真正的非退化谱校准点。于是，局部可展域与全局代数锚定在同一条正实谱证书链上闭合。

\endinput


\endinput
